% The next line makes it so it compiles the main file if I hit the compile button
% when editing this file in TeXworks.
% !TEX root = ../AIDMA.tex
\chapter{Logic}\label{chapter:logicEtc}

\section{Propositional Logic}\label{section:logic}
\subsection{Basic Definitions}\label{section:logicDefs}
\begin{df}
\index{proposition}
\index{boolean!proposition}
\index{truth value}
A {\bf boolean proposition} (or simply {\bf proposition}) is a statement which is either 
{\bf true} or {\bf false} (sometimes abbreviated as {\bf T} or {\bf F}).
We call this the {\bf truth value} of the proposition.
\end{df}
Whether the statement is {\em obviously} true or false does not
enter into the definition. One only needs to know that its certainty
can be established.
\begin{exa}
The following are propositions and their truth values, if known:
\begin{enumerate}[(a)]
\item $7^2 = 49$. ({\bf true}) \item $5>6$. ({\bf false})  
\item If $p$
is a prime then $p$ is odd. ({\bf false}) 
\item There exists infinitely
many primes which are the sum of a square and $1$. (unknown)
%\item There is a God. (unknown)
%\item There is a dog. ({\bf true}) 
\item Dr. Cusack is the Pope. ({\bf false}) 
%\item Every prime that leaves remainder $1$
%when divided by $4$ is the sum of two squares. ({\bf true}) 
\item Every
even integer greater than $6$ is the sum of two distinct primes.
(unknown)
\end{enumerate}
\end{exa}


\begin{rem}
Next you will see the first of many {\bf Exercises}.
These give you an opportunity to solve a problem from start to finish 
and then check your answer with the solution provided.
It is important that you try each of these on your own before looking at the solution.
You will not get as much out of the book if you skip these or jump straight to the answer
without trying them yourself.
\end{rem}

\begin{exer}
Give the truth value of each of the following statements.
\begin{enumerate}[(a)]
\item \blankline{.75} $0=1$.
\item \blankline{.75} 17 is an integer.
%\item \blankline{.75} ``Psych'' is a TV show that aired on the {\em USA network}.
\item \blankline{.75} In 1999, it was possible to buy a red Swingline stapler.
\end{enumerate}
\begin{answer}
(a) false. (b) true. 
%(c) true. 
(c) false.  If you don't know the story behind this, Google it.
\end{answer}
\end{exer}

\begin{rem}
Did you notice the $\star$ in the heading of the previous example?  
This indicates that a solution is provided.
If you are reading the PDF file, clicking on the $\star$ will take you to the solution.
Clicking on the number in the solution will take you back. 

If you are reading the PDF, go to the back of the book to find the solutions.
\end{rem}

\begin{exa}
The following are not propositions, since it is impossible
to assign a {\bf true} or {\bf false} value to them.
\begin{enumerate}[(a)]
\item Whenever I shampoo my camel. 
\item Sit on a potato pan, Otis! 
\item What I am is what I am, are you what you are or what?
\item $x = x + 1$. 
\item This sentence is false.
\end{enumerate}
\end{exa}

\begin{exer}
For each of the following statements, state whether it is true, false, or not a proposition.\\[-.25in]
\begin{enumerate}[(a)]
\item \blankline{.75} i can has cheezburger?
\item \blankline{.75} ``Psych'' was one of the best shows on TV when it was on the air.
\item \blankline{.75} I know, right?
\item \blankline{.75} This is a proposition.
\item \blankline{.75} This is not a proposition.
\end{enumerate}
\begin{answer}
(a) Not a proposition.
(b) I would like to think this is true.  
However, this is not a proposition since not everyone agrees with me.
(c) Also not a proposition.  %(Wait, what?) % Don't include this--they probably won't get it.
(d) true.
(e) false.  This one is a bit tricky to think about it, so the next example will ask you to prove it.
\end{answer}
\end{exer}

\subsection{Compound Propositions}

\begin{df}
\index{operator!logical}
\index{logical operator}
\index{logical!operator}
\index{operator!logical|see{logical, operator}}
\index{proposition!compound}
\index{compound proposition}
A {\bf logical operator} is used to combine one or more
propositions to form a new one.  A proposition formed in this way is called
a {\bf compound proposition}.  
We call the propositions used to form a compound proposition {\bf variables}
for reasons that should become evident shortly.
\end{df}

Next we will discuss the most common logical operators.
Because one of the applications of logic we will be concerned
about is algorithms and programming, when we introduce
notation we will also mention equivalent notations used in 
several common programming languages.
%Some of these will be familiar to you.  When you learned about
%Boolean expressions in your programming courses, you probably saw 
%{\em NOT} (e.g. \verb+if( !list.isEmpty() )+),
%{\em OR} (e.g. \verb+if( x>0 || y>0 )+),
%and 
%{\em AND} (e.g. \verb+if( list.size() > 0 && list.get(0) > 1 )+).
%The notation we use will be different, however.  This is because the symbols you
%are familiar with are specific choices made by whoever created the programming 
%language(s) you learned.  Here we will use standard mathematical notation for 
%the operators.
For each of the following definitions, assume $p$ and $q$ are propositions.

\begin{df}
\index{negation}
\index{logical!operator!negation}
\index{a logicn@$\neg$ (NOT)}
The {\bf negation} (or {\bf NOT}) of $p$, denoted by 
%{\boldmath$\neg p$} 
% Can't get rid of the warning on this one yet.
% Can't find a way around using boldmath or boldsymbol for the neg.
$\boldsymbol{\neg}${\bf\em p} 
is the proposition ``{\bf it is not the case that \boldmath$p$}''.  
$\neg p$ is true when $p$ is false, and vice-versa.
Other notations include $\overline{p}$, $\sim p$,
and $!p$. 
Many programming languages use the last one.
\end{df}

\begin{exa}
If $p$ is ``$x<0$'', then $\neg p$ is 
``It is not the case that $x<0$,'' or ``$x\geq 0$.''
\end{exa}

\begin{rem}
The next example is the first of many {\bf Fill in the details} exercises in 
which {\bf you} need to supply some of the details.   
After you have filled in the blanks, compare your answers with the solutions.  
The answers are often given with semicolons (;) separating the blanks.  
\end{rem}

\vspace*{-.1in} % Stupid spacing trick.

\begin{fib}
Let $p$ be the proposition ``I am learning discrete mathematics.''
Then $\neg p$ is the proposition \blanklinefill.

\noindent
The truth value of $\neg p$ is \blankline{1}.
\begin{answer}
``I am not learning discrete mathematics.''  
You could also have ``It is not the case that I am learning discrete mathematics,''
although it is better to smooth out the English when possible.;
False.  Since you are currently reading the book, you {\em are} learning discrete mathematics.
\end{answer}
\end{fib}

%% Deleted.
%\begin{exer}
%Consider the statement ``This is not a proposition.''
%\begin{enumerate}[(a)]
%\item
%Use the fact that ``This is a proposition'' is a proposition to prove that 
%``This is not a proposition'' is a proposition.  Then prove that its truth value is false.
%
%\pflinethree
%\item
%Use a contradiction proof to prove that ``This is not a proposition'' is a proposition.
%Then prove that its truth value is false.
%
%\pflinethree
%\end{enumerate}
%\begin{answer}
%(a) The statement ``This is not a proposition'' is just the negation of the proposition
%``This is a proposition,'' so it is a proposition.  Since ``This is a proposition'' is true,
%``This is not a proposition'' is false.
%
%(b) Assume that ``This is not a proposition'' is not a proposition.  Then the statement
%``This is not a proposition'' is true, which means it is a proposition.  But we assumed it wasn't
%a proposition.  Since we have a contradiction, our assumption that ``This is not a proposition''
%is not a proposition was incorrect, so it is a proposition.   
%Since it is a proposition that says is isn't, its truth value is false.
%\end{answer}
%\end{exer}

%\newpage

% % Moved to algorithms.
%\begin{exer}
%You need a program to execute some code only if the size of a list is not $0$.
%The variable is named \verb+list+, and its size is \verb+list.size()+.
%Give the expression that should go in the \verb+if+ statement.  
%In fact, give two different expressions that will work.
%\begin{enumerate}
%\item \blanklinefill
%\item \blanklinefill
%\end{enumerate}
%\begin{answer}
%\verb+list.size()!=0+ and \verb+!(list.size()==0)+ are the most obvious solutions.
%\end{answer}
%\end{exer}

\begin{df}
\index{AND}
\index{a logicand@$\wedge$ (AND)}
\index{conjunction}
\index{logical!operator!conjunction}
\index{logical!operator!AND}
The {\bf conjunction} (or {\bf AND}) of $p$ and $q$,
denoted by \boldmath$p\wedge q$, is the proposition ``\boldmath$p$ {\bf and} \boldmath$q$''. 
The conjunction of $p$ and $q$
is true when $p$ and $q$ are both true and false otherwise.
Many programming languages use \verb+&&+ for conjunction.
\end{df}

\begin{exa}
Let $p$ be the proposition ``$x>0$'' and $q$ be the proposition ``$x<10$.''
Then $p\wedge q$ is the proposition ``$x>0$ and $x<10$,'' or ``$0<x<10.$''
%In a Java/C/C++ program, it would be ``\verb+x>0 && x<10+.''
\end{exa}


\vspace*{-.1in} % Stupid spacing trick.

\begin{exa}
Let $p$ be the proposition ``$x<0$'' and $q$ be the proposition ``$x>10$.''
Then $p\wedge q$ is the proposition ``$x<0$ and $x>10$.''

Notice that $p\wedge q$ is always false since if $x<0$, clearly $x\not>10$.
But don't confuse the {\em proposition} with its {\em truth value}.
When you see the statement `$p\wedge q$ is ``$x<0$ and $x>10$'''
and `$p\wedge q$ is false,' these are saying two different things.
The first one is telling us what the proposition is.  
The second one is telling us its truth value.  
`$p\wedge q$ is false' is just a shorthand for saying `$p\wedge q$ has truth value false.'
\end{exa}

\vspace*{-.1in} % Stupid spacing trick.

\begin{fib}
If $p$ is the proposition ``I like cake,'' and $q$ is the proposition ``I like ice cream,''
then $p\wedge q$ is the proposition \blanklinefill.
\begin{answer}
Either 
``I like cake and I like ice cream,'' or
``I like cake and ice cream'' are correct.
\end{answer}
\end{fib}

%% Moved to algorithms
%\begin{exa}
%Write a code fragment that determines whether or not three numbers can
%be the lengths of the sides of a triangle.
%\begin{sol}
%Let $a$, $b$, and $c$ be the numbers.  For simplicity, let's assume they are integers.
%First we must have $a > 0$, $b>0$, and $c>0$. 
%Also, the sum of any two of them must be larger than the third in order to form a triangle.
%More specifically, we need $ a+b>c$, $ b + c>a$, and $c+a>b$.  
%Since we need all of these to be true, this leads to the
%following algorithm.
%\begin{code}
%IsItATriangle(int a, int b, int c) {
%    if(a>0 && b>0 && c>0 && a+b>c && b+c>a && a+c>b) {
%        return true;
%    } else { return false; }
%}\end{code}
%\end{sol}
%\end{exa}

\vspace*{-.1in} % Stupid spacing trick.

\begin{df}
\index{a logicor@$\vee$ (OR)}
\index{OR}
\index{disjunction}
\index{inclusive or}
\index{logical!operator!inclusive or}
\index{logical!operator!OR}
\index{logical!operator!disjunction}
The {\bf disjunction} (or {\bf OR}) of $p$ and $q$,
denoted by \boldmath$p\vee q$, is the proposition ``\boldmath$p$ {\bf or} \boldmath$q$''. 
The disjunction of $p$ and $q$
is false when both $p$ and $q$ are false and true otherwise.  Put another way, if $p$ is true, $q$ is true,
or both are true, the disjunction is true.
Many programming languages use \verb+||+ for disjunction.
\end{df}


\begin{exa}
Let $p$ be the proposition ``$x<5$'' and $q$ be the proposition ``$x>15$.''
Then $p\vee q$ is the proposition ``$x<5$ or $x>15$.''
In a Java/C/C++ program, it would be ``\verb+x<5 || x>15+.''
\end{exa}

\vspace*{-.1in} % Stupid spacing trick.

\begin{fib}
Let $p$ be the proposition ``$x>0$'' and $q$ be the proposition ``$x<10$.''
Then $p\vee q$ is the proposition \blanklinefill.\\
Notice that $p\vee q$ is always \blankline{1} since 
it is \blankline{1} if $x>0$, and if $x\not>0$, then clearly \blankline{1.5}, so it is \blankline{1}
then as well.
\begin{answer}
``$x>0$ or $x<10$''; true; true; $x<10$; true.
\end{answer}
\end{fib}

\vspace*{-.1in} % Stupid spacing trick.

%
%\begin{exer}
%Let $p$ be ``you must be at least 48 inches tall to ride the roller coaster,'' and
%$q$ be ``you must be at least 18 years old to ride the roller coaster.''
%Express each of the following propositions in English. 
%\begin{enumerate}
%\item
%$\neg p$ is \blanklinefill\\
%\blanklinefill\\
%\blanklinefill
%\item
%$p\vee q$ is \blanklinefill\\
%\blanklinefill\\
%\blanklinefill
%\item
%$p\wedge q$ is \blanklinefill\\
%\blanklinefill\\
%\blanklinefill
%\end{enumerate}
%\begin{answer}
%(a) ``You do not have to be at least 48 inches tall to ride the roller coaster.''
%It is {\em not} ``You must be at most 48 inches tall to ride the roller coaster.''
%(b) ``You must be at least 48 inches tall or 18 years old to ride the roller coaster.''
%(c) ``You must be at least 48 inches tall and at least 18 years old to ride the roller coaster.''
%\end{answer}
%\end{exer}

\begin{exer}
Let $p$ be ``Jill is tall,'' and
$q$ be ``Jill is smart.''
Express each of the following propositions in English. 
\begin{enumerate}[(a)]
\item
$\neg p$ is \blanklinefill
\item
$p\vee q$ is \blanklinefill
\item
$p\wedge q$ is \blanklinefill
\item
$p\wedge \neg q$ is \blanklinefill
\item
$\neg(p\wedge q)$ is \blanklinefill
\end{enumerate}
\begin{answer}
(a) ``It is not the case that Jill is tall,'' 
which is awkward, so could be shorted to
``Jill is not tall'' or perhaps ``Jill is short.''
(b) ``Jill is tall or Jill is smart,'' or more compactly, 
``Jill is tall or smart.''
(c) ``Jill is tall and Jill is smart,'' or more compactly, 
``Jill is tall and smart.''
(d) ``Jill is tall and Jill is not smart.''
(e) This one is a bit tricky, and we will see a tool shortly 
that will make it easier. But for now, hopefully you can see
that is would be ``it is not the case that Jill is tall and Jill is smart.''
Notice that an {\em incorrect} answer would be
``Jill is neither tall nor smart.'' We will see why later.
\end{answer}
\end{exer}

\vspace*{-.1in} % Stupid spacing trick.

%% MOved to algorithms
%\begin{exer}
%Give an algorithm that will return true if an array of integers either starts or ends with a $0$, and
%false otherwise.  Assume array indexing starts at $0$ and that the array is of length $n$.
%Use only one conditional statement.  Be sure to deal with the possibility of an empty array.
%\begin{code}
%    boolean startsOrEndsWithZero(int[] a, int n) {
%
%
%
%
%
%
%
%    }
%\end{code}
%\begin{answer}
%Here is one possible solution.  Note that the parentheses are necessary.
%\begin{code}
%    boolean startsOrEndsWithZero(int[] a, int n) {
%        if( n>0 && (a[0]==0 || a[n-1]==0) ) {
%            return true;
%        else {
%            return false;
%        }
%    }
%\end{code}
%\end{answer}
%\end{exer}
%
%\begin{ques}
%Does the solution given for the previous exercise properly deal with arrays of size $0$ and $1$?  
%Prove it.\\
%\answerlinefour
%\begin{answer}
%The solution uses the expression  \verb+n>0 && (a[0]==0 || a[n-1]==0)+.
%If $n=0$, the expression is false because of the \verb+&&+, so the algorithm returns false as it should
%since an array with no elements certainly does not begin or end with a $0$.
%If $n=1$, first note that $n-1=0$, so $a[0]$ and $a[n-1]$ refer to the same element.
%Although this is redundant, it isn't a problem.  
%If $a[0]=0$, the expression evaluates to $T\wedge (T\vee T)=T$, and the algorithm returns
%true as expected.  
%If $a[0]\neq 0$, the expression evaluates to $T\wedge (F\vee F)=F$, and the algorithm returns
%false as expected.
%\end{answer}
%\end{ques}

\begin{df}
\index{a logicxor@$\oplus$ (XOR)}
\index{exclusive or}
\index{XOR}
\index{logical!operator!exclusive or}
\index{logical!operator!XOR}
The {\bf exclusive or} (or {\bf XOR}) of $p$ and $q$,
denoted by \boldmath$p\oplus q$, is the proposition ``{\bf \boldmath$p$ is true or \boldmath$q$ is true, but not both}''.
The exclusive or of $p$ and $q$ is true when exactly one of $p$ or $q$ is true.
Put another way, the exclusive or of $p$ and $q$ is true iff $p$ and $q$ have different truth values.
\end{df}

\vspace*{-.1in} % Stupid spacing trick.

\begin{exa}
Let $p$ be the proposition ``$x>10$'' and $q$ be the proposition ``$x<20$.''
Then $p\oplus q$ is the proposition ``$x>10$ or $x<20$, but not both.''
\end{exa}

\vspace*{-.1in} % Stupid spacing trick.

\begin{rem}
Notice that $\vee$ is an {\bf inclusive or}, meaning that it is true if both are true, 
whereas  $\oplus$ is an {\bf exclusive or}, meaning it is false if both are true.
The difference between $\vee$ and $\oplus$ is complicated by the fact that in English,
 the word ``or'' to can mean either of these depending on context.  For instance, if your mother
tells you ``you can have cake or ice cream'' she is likely using the exclusive or, 
whereas a prerequisite of ``Math 110 or demonstrated competency with algebra'' clearly
has the inclusive or in mind.
\end{rem}


\begin{exer}
For each of the following, is the {\em or} inclusive or exclusive?  
Answer {\bf OR} or {\bf XOR} for each.
\begin{enumerate}[(a)]
\item \blankline{1} The special includes your choice of a salad or fries.
\item \blankline{1} The list is empty or the first element is zero.
\item \blankline{1} The first list is empty or the second list is empty.
\item \blankline{1} You need to take probability or statistics before taking this class.
\item \blankline{1} You can get credit for either Math 111 or Math 222.
\end{enumerate}
\begin{answer}
(a) XOR; 
(b) OR.  This one is a little tricky because parts can't be simultaneously true
so it sounds like an XOR.  But since the point of the statement is not to {\em prevent}
both from being true, it is an OR.
(c) Without more context, this one is difficult to answer.  I would suspect that most
of the time this is probably OR.  The purpose of this example is to demonstrate that
sometimes life contains ambiguities.  This is particularly true with software specifications.
Generally speaking, you should {\em not} just assume one of the alternative possibilities.
Instead, get the ambiguity clarified.
(d) When course prerequisites are involved, OR is almost certainly in mind.
(e) The way this is phrased, it is almost certainly an XOR.
\end{answer}
\end{exer}



\begin{exer}
Let $p$ be ``list 1 is empty'' and $q$ be ``list 2 is empty.''
Explain the difference in meaning between $p\vee q$ and $p\oplus q$.\\
\answerlinethree
\begin{answer}
$p\vee q$ is ``either list 1 or list 2 is empty.''  To be completely unambiguous, 
you could rephrase it as ``at least one of list 1 or list 2 is empty.''
$p\oplus q$ is ``either list 1 or list 2 is empty, but not both,'' or
``precisely one of list 1 or list 2 is empty.''
They are different because if both lists are empty, $p\vee q$ is true, but $p\oplus q$ is false.
\end{answer}
\end{exer}

\begin{rem}
The {\bf Question} examples are similar to the {\bf Evaluate} ones except that they ask a specific question.
Write down your answer in the space provided and then compare your answer with the one in the solutions.
\end{rem}

\begin{ques}
Let $p$ be the proposition ``$x<5$'' and $q$ be the proposition ``$x>15$.''
\begin{enumerate}[(a)]
\item
Do the statements $p\vee q$ and $p\oplus q$ mean the same thing? Explain.\\
\answerlinetwo
\item
Practically speaking, are $p\vee q$ and $p\oplus q$ the same?  Explain.\\
\answerlinetwo
\end{enumerate}
\begin{answer}
(a) No.  If $p$ and $q$ are both true, then $p\vee q$ is true, but $p\oplus q$ is false, so they do not mean the same thing.
(b) We have to be very careful here.  In general, the answer to this would be absolutely not
(we'll discuss this more next).
However, for {\em this particular $p$ and $q$}, they actually essentially are the same.
But the reason is that it is impossible for $x$ to be less than $5$ and greater than $15$ at
the same time.  In other words, $p$ and $q$ can't both be true at the same time.
The only other way for $p\oplus q$ to be false is if both $p$ and $q$ are false, which is 
exactly when $p\vee q$ is false.
\end{answer}
\end{ques}

\newpage

XOR is not used as often as AND and OR in logical expressions in programs.  
Some languages have an XOR operator and some do not.  
The issue gets blurry because some languages, like Java, have an explicit Boolean type, 
while others, like C and C++, do not.
All of these languages have a {\em bitwise XOR} operator, but this is not the same thing as a
{\em logical XOR} operator. We will return to this topic later.
In the next section we will see how to implement $\oplus$ using $\vee$, $\wedge$, and $\neg$.

\begin{df}
\index{a logicyy@$\rightarrow$ (conditional)}
\index{conditional statement}
\index{logical!operator!conditional}
The {\bf conditional statement} (or {\bf implies} or {\bf implication}) involving $p$ and $q$,
denoted by \boldmath$p \rightarrow q$, is the proposition ``{\bf if \boldmath$p$, then \boldmath$q$}''.
It is false when $p$ is true and $q$ is false, and true otherwise.
In the statement $p \rightarrow q$, we call $p$ the {\bf premise} (or {\bf hypothesis} or {\bf antecedent})
and $q$ the {\bf conclusion} (or {\bf consequence}).
\end{df}

\begin{exa}
Let $p$ be ``you earn at least 94\%,'' and $q$ be ``you will receive an A.''
Then $p\rightarrow q$ is the proposition ``If you earn at least 94\%, then you will receive an A.''
\end{exa}

It is important to realize that $p\rightarrow q$ and $q\rightarrow p$ are not always equivalent.

\begin{exa}
Let $p$ be ``you earn at least 94\%,'' 
and $q$ be ``you will receive an A.''
Then $p\rightarrow q$ is the proposition ``If you earn at least 94\%, then you will receive an A,'' and 
$q\rightarrow p$ is the proposition ``If you receive an A, then you earned at least 94\%.'' Although they may sound equivalent, they are not. 
Consider the possibility that it is true that receiving at least 93\% results in an A. Then $p\rightarrow q$ is true, but $q\rightarrow p$ is false.
\end{exa}

\begin{ques}
Assume that the proposition ``If you earn at least 94\% in this class, then
you will receive an A'' is true.
\begin{enumerate}[(a)]
\item What grade will you get if you earn 94\%?  Explain.\\
\answerlinetwo
\item If you receive an A, did you earn at least 94\%? Explain.\\
\answerlinetwo
\item If you don't earn 94\%, does that mean you didn't get an A? Explain.\\
\answerlinetwo
\end{enumerate}
\begin{answer}
(a) An A. This is pretty obvious since we earned 94\% and if we
earn 94\%, then we will get an A.
(b) We can't be sure.  We know that earning 94\% is enough for an A, but we don't know
whether or not there are other ways of earning an A.
(c) We can't be sure.  If the premise is false, we don't know anything about conclusion.
\end{answer}
\end{ques}

\begin{exa}
Translating between an English sentence and a mathematical expression
can sometimes be tricky with conditional statements.
Again, let 
$p$ be ``you earn at least 94\%,'' and $q$ be ``you will receive an A.''
Then the sentence ``You will receive an A whenever you earn at least 94\%''
is $p\rightarrow q$, and not $q\rightarrow p$ since it is expressing
the same idea as the sentence ``If you earn at least 94\%, you will receive an A.''
\end{exa}

\begin{rem}
The {\bf conditional} statement is by far the one that is the most difficult to get a handle on for at least
two reasons.  First, the conditional statement $p \rightarrow q$ is not saying anything about $p$ or $q$ by themselves.
It is only saying that if $p$ is true, then $q$ has to also be true.  
It doesn't say anything about the case that $p$ is
not true.  This brings us to the second reason.
Should $F \rightarrow T$ be true or false?  Although it seems counterintuitive to some, it should be true.  
Again, $p \rightarrow q$ is telling us about the value of $q$ when $p$ is true 
(i.e., if $p$ is true, the $q$ {\em must be} true).  
What does it tell us if $p$ is false?  Nothing.  As strange as it might seem, when $p$ is
false, the whole statement is true regardless of the truth value of $q$.

If you are still confused, you can simply fall back on the formal definition:
{\bf \boldmath$p \rightarrow q$ is false when \boldmath$p$ is true and \boldmath$q$ is false, 
and is true otherwise.}
In other words, if interpreting $p \rightarrow q$ as 
the English sentence ``$p$ implies $q$'' is more harmful than helpful in understanding
the concept, don't worry about why it doesn't make sense and just remember 
the definition.\footnote{In mathematics, 
terms are usually chosen so they make sense immediately.  
Sometimes this is not possible (if the concept is very complicated or it 
doesn't relate to anything familiar). Sometimes a term
is poorly defined but the definition sticks because of prior use.  
Sometimes it makes sense to some people and not to others, 
probably based on a person's background.
I think this last possibility may be the reason in this case.}
\end{rem}

We will learn more about the conditional statements and
statements related to it in the chapter on proofs where it
is particularly relevant. 

\begin{df}
\index{a logiczz@$\leftrightarrow$ (biconditional)}
\index{biconditional}
\index{logical!operator!biconditional}
The {\bf biconditional statement}  involving $p$ and $q$,
denoted by \boldmath$p \leftrightarrow q$, is the proposition ``{\bf \boldmath$p$ if and only if \boldmath$q$}''
(or abbreviated as ``{\bf \boldmath$p$ iff \boldmath$q$}'').
It is true when $p$ and $q$ have the same truth value, and false otherwise.
\end{df}

\begin{exa}
Let $p$ be ``you earn at least 94\%,'' and $q$ be ``you receive an A.''
Then $p\leftrightarrow q$ is the proposition ``You earn at least 94\% if and only if you receive an A.''
\end{exa}

\begin{ques}
Assume that the proposition ``You will receive an A
if and only if you earn at least 94\%'' is true.
\begin{enumerate}[(a)]
\item What grade will you get if you earn 94\%?\\
\answerline
\item If you receive an A, did you earn at least 94\%?\\
\answerline
\item If you don't earn at least 94\%, does that mean you didn't get an A?\\
\answerline
\end{enumerate}
\begin{answer}
(a) An A.
(b) Yes. Because it is a biconditional statement that we assumed to be true, 
the statements ``you will receive an A in the course'' and ``you earn at least 94\%'' have the
same truth value.  Since the former is true, the latter has to be true.
(c) Yes.  Notice that $p\leftrightarrow q$ is equivalent to $\neg p \leftrightarrow \neg q$
(You should convince yourself that this is true).  
Thus the statements ``you don't earn at least 94\%'' and ``you didn't get an A'' have the same
truth value.    
\end{answer}
\end{ques}

Now let's bring all of these operations together with a few more examples.

\begin{exa}
Let $a$ be the proposition ``I will eat my socks,'' $b$ be ``It is snowing,'' and
$c$ be ``I will go jogging.''
Here are some compound propositions involving $a$, $b$, and $c$, written using these
variables and operators and in English.\\
\begin{tabular}{|l|p{.63\textwidth}|}
\hline
With Variables/Operators&In English\\
\hline
\hline
$(b \vee \neg b)  \rightarrow c$ &
Whether or not it is snowing, I will go jogging.\\
\hline
$b  \rightarrow \neg c$ & 
If it is snowing, I will not go jogging.\\
\hline
$b  \rightarrow (a \wedge \neg c)$&
If it is snowing, I will eat my socks, but I will not go jogging.\\
\hline
$a \leftrightarrow c$&
When I eat my socks I go jogging, and when I go jogging I eat my socks.\\
&{\em or} I eat my socks if and only if I go jogging.\\
\hline
\end{tabular}
\end{exa}

\begin{fib}
Let $p$ be the proposition ``{\em Iron Man} is on TV,'' $q$ be ``I will watch {\em Iron Man},'' and
$r$ be ``I own {\em Iron Man} on DVD.''
Fill in the missing information in the following table.\\
\begin{tabular}{|l|p{.63\textwidth}|}
\hline
With Variables/Operators&In English\\
\hline
\hline
$p\rightarrow q$&\\[.35in]
\hline
&If I don't own {\em Iron Man} on DVD and it is on TV, I will watch it.\\[.25in]
\hline
$p\wedge r \wedge\neg q$&\\[.35in]
\hline
&I will watch {\em Iron Man} every time it is on TV, and that is the only time I watch it.\\[.25in]
\hline&I will watch {\em Iron Man} if I own the DVD.\\[.25in]
\hline
\end{tabular}
\begin{answer}
The answers are in {\bf bold}.\\
\begin{tabular}[t]{|l|p{.575\textwidth}|}
\hline
With Variables/Operators&In English\\
\hline
\hline
$p\rightarrow q$&{\bf If {\em Iron Man} is on TV, then I will watch it.}\\
\hline
\boldmath$(\neg r\wedge p)\rightarrow q$&If I don't own {\em Iron Man} on DVD and it is on TV, I will watch it.\\
\hline
$p\wedge r \wedge\neg q$&{\bf {\em Iron Man} is on TV and I own the DVD, but I won't watch it.}\\
\hline
\boldmath$q\leftrightarrow p$&I will watch {\em Iron Man} every time it is on TV, and that is the only time I watch it.\\
\hline
\boldmath$r\rightarrow q$&I will watch {\em Iron Man} if I own the DVD.\\
\hline
\end{tabular}
\end{answer}
\end{fib}


\subsection{Truth Tables}

Sometimes we will find it useful to think of compound propositions in terms of {\em truth tables}.

\begin{df}
\index{truth table}
A {\bf truth table} is a table that shows the truth value of a 
compound proposition for all possible combinations of
truth assignments to the variables in the proposition. If there are $n$
variables, the truth table will have $2^n$ rows.
\end{df}

\begin{table}[h]
\begin{minipage}[t]{.7\textwidth}
\begin{tabular}[t]{p{\textwidth}}
The truth table for $\neg$ is given in Table~\ref{tab:not} and
the truth tables for all of the other operators we just defined are given in Table~\ref{tab:2vars}.
In the latter table, the first two columns are the possible values of the two variables, and the last 5 columns are
the values for each of the two-variable compound propositions we just defined for the given inputs.
\end{tabular}
\end{minipage}
\hfill
\begin{minipage}[t]{.2\textwidth}
\begin{tabular}[t]{|c|c|}
\hline $p$ & $\neg p$\\
\hline 
$T$ & $F$\\
$F$ & $T$\\
\hline
\end{tabular}
\footnotesize\caption{\\Truth table for $\neg$}\label{tab:not}
\end{minipage}
\end{table}

\begin{table}[h]
\centering
\begin{tabular}{|cc||c|c|c|c|c|}
\hline $p$ & $q$ & $(p \wedge q)$ & $(p \vee q)$ &
$p\oplus q$&
$(p  \rightarrow q)$ & $(p  \leftrightarrow q)$ \\
\hline 
$T$ & $T$& $T$ &$T$&$F$&$T$&$T$ \\
$T$ & $F$& $F$ &$T$&$T$&$F$&$F$ \\
$F$ & $T$& $F$ &$T$&$T$&$T$&$F$ \\
$F$ & $F$& $F$ &$F$&$F$&$T$&$T$ \\
\hline
\end{tabular}
\footnotesize\caption{Truth tables for the two-variable operators}\label{tab:2vars}
\end{table}

\begin{exa}\label{exa:truth_table_example}
Construct the truth table of the proposition $a \vee  (\neg b \wedge c)$.
\begin{sol}
Since there are three variables, the truth table will have
$2^3 = 8$ rows. Here is the truth table, with several helpful intermediate columns.

\centering
\begin{tabular}{|ccc||c|c|c|}
\hline $a$ & $b$ & $c$ & $\neg b$ & $\neg b \wedge c$ & $a \vee
(\neg b
\wedge c)$ \\
\hline
$T$ & $T$ & $T$ & $F$ & $F $ & $T $ \\
$T$ & $T$ & $F$ & $F $ & $F $ & $T $ \\
$T$ & $F$ & $T$ & $T $ & $T $ & $T $ \\
$T$ & $F$ & $F$ & $ T$ & $F $ & $T $ \\
$F$ & $T$ & $T$ & $F $ & $F $ & $F $ \\
$F$ & $T$ & $F$ & $F $ & $F $ & $F $ \\
$F$ & $F$ & $T$ & $T $ & $T $ & $T $ \\
$F$ & $F$ & $F$ & $T $ & $F $ & $F $ \\
\hline
\end{tabular}
\end{sol}
\end{exa}


\begin{rem}
Notice that there are several columns in the truth table besides the columns for
the variables and the column for the proposition we are interested in.
These are ``helper'' or ``intermediate'' columns (those are not official definitions).
Their purpose is simply to help us compute the final column more easily and without
(hopefully) making any mistakes.
\end{rem}


\begin{exer}
Construct the truth table for  $(p\rightarrow q) \wedge q$. 

{\large
$$ 
\begin{array}{|cc||c|c|} 
\hline 
p & q &p\rightarrow q &(p \rightarrow q) \wedge q\\
\hline
\hline
T & T &  &  \\
\hline
T & F &  &  \\
\hline
F & T &  &  \\
\hline
F & F &  &  \\
\hline
\end{array}
$$
}
\begin{answer}
Here is the truth table with one (optional) intermediate column.
$$ \begin{array}{|cc||c|c|} \hline p & q & p  \rightarrow q & (p \rightarrow q) \wedge q \\
\hline
T & T & T & T\\
T & F & F & F \\
F & T & T & T \\
F & F & T & F \\
\hline
\end{array}$$
\end{answer}
\end{exer}

\begin{rem}
As long as all possible values of the variables are included,
the order of the rows of a truth table does not matter. 
However,  by convention one of two orderings is usually used.  
Since there is an interesting connection to the binary representation of numbers, 
let's take a closer look at this connection in the next example.
\end{rem}

\begin{exa}[Ordering the rows of a Truth Table]\label{exa:rowOrder}
Notice that the values of the variables can be used to construct an index for each row.  
%So if proposition involves two variables, 
%the values in the first two columns are used as a sort of index.
We can do this by thinking of each T as a $1$ 
and each F as a $0$, and treating the columns as a binary number.
The rows will then be listed either in order or (more commonly) in 
reverse order.
%We can order the rows by assigning a number to each row based on the values in these columns.
%The order used here essentially computes an index as follows:  For the ``index'' columns, think of each T as a $1$ 
%and each F as a $0$.  Now treat the numbers in these columns as binary numbers 
%and order the rows appropriately.  
For instance, if there are three variables, we can think of it as shown
in the following table.

\begin{center}
\begin{tabular}{|ccc||ccc|c|}
\hline $a$ & $b$ & $c$ & &  &&index \\
\hline
$T$ & $T$ & $T$ &1&1&1&7\\
$T$ & $T$ & $F$ &1&1&0&6\\
$T$ & $F$ & $T$ &1&0&1&5\\
$T$ & $F$ & $F$ &1&0&0&4\\
$F$ & $T$ & $T$ &0&1&1&3\\
$F$ & $T$ & $F$ &0&1&0&2\\
$F$ & $F$ & $T$ &0&0&1&1\\
$F$ & $F$ & $F$ &0&0&0&0\\
%$T$ & $T$ & $T$ &0 & 0&0&0\\
%$T$ & $T$ & $F$ &0&0 &1&1\\
%$T$ & $F$ & $T$ &0&1&0&2\\
%$T$ & $F$ & $F$ &0&1 &1 &3\\
%$F$ & $T$ & $T$ & 1&0&0&4\\
%$F$ & $T$ & $F$ &1&0 &1&5\\
%$F$ & $F$ & $T$ &1 &1 &0&6\\
%$F$ & $F$ & $F$ &1& 1&1&7\\
\hline
\end{tabular}
\end{center}
{\em This is the ordering you should follow so that you can easily check your answers with those in
the solutions.}  
It also makes grading your homework easier.

%The other common ordering does the same thing, but maps T to $1$  and F to $0$.  

There is also a way of thinking about this recursively. 
Given an ordering for a table with
$n$ variables, we can compute an ordering for a table with $n+1$ variables as follows.
Make two copies of the columns corresponding to the $n$ variables, 
appending a T to the beginning of the
first copy, and an F to the beginning of the second copy.
\end{exa}



\begin{exer}\label{exa:truth_table_exampleb}
Construct the truth table of the proposition $(a \vee  \neg b) \wedge c$.
You're on your own this time to supply all of the details.

\vspace*{3in}
\begin{answer}
Here is the truth table with two intermediate columns.
For consistency, your table should have the rows in the same order.\\
\begin{tabular}[t]{|ccc||c|c|c|}
\hline $a$ & $b$ & $c$ & $\neg b$ & $a\vee \neg b$ & $(a \vee\neg b)\wedge c)$ \\
\hline
$T$ & $T$ & $T$ & $F $ & $T $ & $T $ \\
$T$ & $T$ & $F$ & $F $ & $T $ & $F $ \\
$T$ & $F$ & $T$ & $T $ & $T $ & $T $ \\
$T$ & $F$ & $F$ & $T $ & $T $ & $F $ \\
$F$ & $T$ & $T$ & $F $ & $F $ & $F $ \\
$F$ & $T$ & $F$ & $F $ & $F $ & $F $ \\
$F$ & $F$ & $T$ & $T $ & $T $ & $T $ \\
$F$ & $F$ & $F$ & $T $ & $T $ & $F $ \\
\hline
\end{tabular}
\end{answer}
\end{exer}

\subsection{Precedence Rules}
Consider the compound proposition $a \vee  \neg b \wedge c$.
Should this be interpreted as  $a \vee  (\neg b \wedge c)$, $(a \vee  \neg b) \wedge c$,
or even possibly $a \vee  \neg (b \wedge c)$?  Does it even matter? 
You already know that $3+(4*5)\neq (3+4)*5$, so it should not surprise you
that where you put the parentheses in logical expressions matters, too.
In fact, Example~\ref{exa:truth_table_example} gives the truth table for one
of these and you just computed the truth table for another one in Exercise~\ref{exa:truth_table_exampleb}.
If you compare them, you will see that they are not the same.  The third interpretation is also different
from both of these.

To correctly interpret compound propositions, the operators have an {\em order of precedence}.
The order is $\neg$, $\wedge$, $\oplus$, $\vee$, $\rightarrow$, and $\leftrightarrow$.
Also, $\neg$ has right-to-left associativity, all other operators listed
have left-to-right associativity.\index{precedence, logical operators}
Based on these rules, 
the correct way to interpret $a \vee  \neg b \wedge c$ is $a \vee  ((\neg b) \wedge c)$.

It is important to know the precedence rules for the logical operators 
(or at least be able to look it up) so you can properly interpret complex logical expressions.
However, I generally prefer to always use enough parentheses to make it immediately clear, 
especially when I am writing code. It isn't difficult to remember that $\neg$ is first 
(that is, it always applies to what is immediately after it) so sometimes I don't use parentheses for it.

\begin{exa}
According to the precedence rules, 
$\neg a  \rightarrow a \vee b$ should be interpreted as $(\neg a)  \rightarrow
(a\vee b)$.
\end{exa} 

\vspace*{-.1in} % Stupid spacing trick.

\begin{exa}
According to the precedence rules, 
$a \wedge \neg b  \rightarrow c$ should be interpreted as  $(a \wedge (\neg b))
 \rightarrow c$.
\end{exa}

\vspace*{-.1in} % Stupid spacing trick.

\begin{exer}
According to the precedence rules, 
how should $a \wedge  b \vee c$  be interpreted?\\
\answerline
\begin{answer}
$(a \wedge  b)\vee c$
\end{answer}
\end{exer}

\vspace*{-.1in} % Stupid spacing trick.

\begin{ques}
Are  $(a\wedge b)\vee c$ and $a \wedge ( b\vee c)$ equivalent?  Prove your answer.\\
\answerlinetwo
\begin{answer}
They are not equivalent.  For instance, when $a=F$, $b=F$, and $c=T$, $(a\wedge b)\vee c$ is true
but $a \wedge ( b\vee c)$ is false.
\end{answer}
\end{ques}

\vspace*{-.1in} % Stupid spacing trick.

\begin{rem}
The next example is an {\bf Evaluate} exercise.  
These exercises give a problem and then provide one or more solutions to the problem 
based on previous student solutions.
Your job is to evaluate each solution by finding any mistakes.  
Mistakes include not only incorrect algebra and logic, but also unclear presentation, 
skipped steps, incorrect assumptions, over-simplification, etc.
When you come across these examples you should write down every error you can find.  Once you are pretty sure you
know all of the problems (if there are any), compare your evaluation to the one given in the solutions.
Note that sometimes the given solutions are correct!
\end{rem}

\vspace*{-.1in} % Stupid spacing trick.

\begin{eval}
According to the associativity rules, how should
$a \rightarrow b  \rightarrow c$  be interpreted?
\begin{ssol}
It should be interpreted as $(a \rightarrow b) \rightarrow c$.
However, $a \rightarrow (b \rightarrow c)$ is equivalent, so it really doesn't matter.
\end{ssol}
\evallinetwo
\begin{answer}
Since $(a \rightarrow b) \rightarrow c$ is how it should be interpreted, the first statement is correct.
The second statement is incorrect.  We'll leave it to you to find true values for $a$, $b$, and $c$
that result in these two parenthesizations having different truth values.
\end{answer}
\end{eval}

\vspace*{-.1in} % Stupid spacing trick.

\newpage
\section{Propositional Equivalence}\label{section:equiv}
We have already informally discussed two propositions being {\em equivalent}.
In this section, we will formally develop the notion of what it means for
two propositions to be {\em equivalent} (or, more formally,  {\em logically equivalent}).
We will also provide you with a list of the most important logical equivalences,
along with some examples of some that aren't necessarily as important, but make interesting examples.
But first, we need some new terminology.

\begin{df}
\index{tautology}
\index{contradiction}
\index{contingency}
A proposition that is always true is called a {\bf tautology}.  
One that is always false is a {\bf contradiction}.  
Finally, one that is neither of these is called  a {\bf contingency}.
\end{df}

\begin{exa}
Assume that $x$ is a real number.
\begin{enumerate}[(a)]
\item
The proposition ``$x<0$'' is a contingency since its truth depends on the value of $x$.
\item
The proposition ``$x^2<0$'' is a contradiction since it is false no matter what $x$ is.
\item
The proposition ``$x^2\geq 0$'' is a tautology since it is true no matter what $x$ is.
\end{enumerate}
\end{exa}

\begin{fib}
State whether each of the following propositions is a tautology, contradiction, or contingency.
Give a brief justification.
\begin{enumerate}[(a)]
\item $p\vee\neg p$ is a \blankline{1} since either $p$ or $\neg p$ has to be true.
\item $p\wedge\neg p$ is a \blankline{1} since \blanklinefill.
\item $p\vee q$ is a \blankline{1} since \blanklinefill.
\end{enumerate}
\begin{answer}
(a) tautology (b) contradiction;  $p$ and $\neg p$ cannot both be true.
(c) contingency; it can be either true or false depending on the truth values of $p$ and $q$.
\end{answer}
\end{fib}

We will cover proofs more formally later, but for now we will
informally introduce two proof techniques involving propositional
logic. 
To prove something is a tautology, one must prove that it is always true.  One way to do this is
to show that the proposition is true for every row of the truth table.  Another way is to
argue (without using a truth table) that the proposition is always true, often using a {\em proof by cases}. 
This is exactly what it sounds like: 
consider every possibility, and show that in all cases we get true.

\begin{exa}
Prove that $p\vee \neg p$ is a tautology.\\
Here are several proofs.
\begin{pfenum}
\item
Since every row in the following truth table for $p\vee \neg p$ is $T$, it is a tautology.
$$\begin{array}{|c||c|c|c|}
\hline
p&\neg p&p\vee \neg p\\
\hline
T&F&T\\
F&T&T\\
\hline
\end{array}$$
\item
By definition of disjunction, if $p$ is true, then $p\vee \neg p$ is true.  On the other hand, if
$p$ is false, $\neg p$ is true.  In this case, $p\vee \neg p$ is still true,
again by definition of disjunction.
Since $p\vee \neg p$ is true regardless of the value of $p$, it is a tautology. 
\end{pfenum}
\end{exa}

\begin{eval}
Prove that $[p\wedge (p  \rightarrow q)] \rightarrow q$ is a tautology.
\begin{spfenum}
\item 
$$\begin{array}{|c|c||c|c|c|}
\hline
p&q&p\rightarrow q&p\wedge(p\rightarrow q)&p\wedge(p\rightarrow q)\rightarrow q\\
\hline
T&T&T&T&T\\
T&F&F&F&T\\
F&T&T&F&T\\
F&F&T&F&T\\
\hline
\end{array}$$
\evallinetwo
\item
One way to show that $p\wedge(p\rightarrow q)\rightarrow q$ is indeed a tautology is by filling out a truth table,
as follows:
$$\begin{array}{|c|c||c|c|c|}
\hline
p&q&p\rightarrow q&p\wedge(p\rightarrow q)&p\wedge(p\rightarrow q)\rightarrow q\\
\hline
T&T&T&T&T\\
T&F&F&F&T\\
F&T&T&F&T\\
F&F&T&F&T\\
\hline
\end{array}$$
Since they all return true for $p\wedge(p\rightarrow q)\rightarrow q$, this proves that it is a tautology.\\
\evallinethree
\item
One way to prove that this is a tautology is to make a couple of assumptions.
First, since we know that for any statement $x\rightarrow y$ where $y$ is true, then $x$ can
be either true or false.  So let us assume that $q$ is false for this case.
From the left side of the statement, if $p$ is true, we would have true and (true implies false), which is false,
thus we would have false implies false, which is true, and if $p$ is false, then we would 
have false and (false implies true), which comes out false.
So in both cases where $q$ is false, the statement equals out to false implies false, which is true,
thus all four cases are true, thereby proving that $p\wedge(p\rightarrow q)\rightarrow q$ is a tautology.\\
\evallinethree
\item
Since an implication can only be false when the premise is true and the conclusion is false, we only need
to prove that this can't happen. So let's assume that $p\wedge(p\rightarrow q)$ is true but that
$q$ is false.  Since $p\wedge(p\rightarrow q)$ is true, $p$ is true and $p\rightarrow q$ is true (by definition of
conjunction).  But if $p$ is true and $q$ is false, $p\rightarrow q$ is false. 
This is a contradiction, so it must be the case that our assumption that 
$p\wedge(p\rightarrow q)$ is true but that $q$ is false is incorrect.
Since that was the only possible way for $p\wedge(p\rightarrow q)\rightarrow q$ to be false, it cannot be false.
Therefore it is a tautology.\\
\evallinethree
\item
Because `merica.\\
\evallinetwo
\end{spfenum}
\begin{answer}
\begin{pfevalenum}
\item
Nice truth table, but what does it {\em mean}?  It is just a bunch of symbols on a page.
Why does this truth table prove that the proposition is a tautology?
The proof needs to include a sentence or two to make the connection between the truth table and
the proposition being a tautology.
\item
 This is mostly correct, but the phrasing could be improved.  For instance, the phrase
`they all return true' is problematic.  Who/what are `they'?  And what does it mean that they
`return' true?   Propositions don't `return' anything.
Replace `Since they all return true' with `Since every row of the table is true' and the proof would be good.
\item
While I applaud the attempt at completeness, this proof is way too complicated.  It is hard to
understand because of the incredibly long sentences and the mathematical statements written in English
in the middle of sentences.  But I suppose that technically speaking it is correct.
Here are a few specific examples of problems with the proof (not exhaustive).  
The first three sentences are confusing as stated.  The point that the 
author is trying to make is that whenever $q$ is true, the statement must be true regardless of the value of $p$, 
so there is nothing further to verify.  Thus the only case left is when $q$ is false.  This point could be made
with far few words and more clearly.
The phrase `we would have true and (true implies false), which is false,' is very confusing, as are a few similar
statements in the proof.  The problem is that the writer is trying to express mathematical statements in 
sentence form instead of using mathematical notation.  
There is a reason we learn mathematical notation--to use it!  
\item
This proof is correct and is not too difficult to understand.  It is a lot better than the previous proof for a few
reasons.  First of all, it starts off in a better place--focusing in on the single case of importance.  
Second, it uses the appropriate mathematical notation and refers to definitions and previous facts to clarify the argument.
\item
While I appreciate the patriotism (in case you don't know, some people use `merica as a shorthand for America),
this has nothing to do with the question.  Sorry, no points for you!
By the way, I did {\em not} make this solution up.  Although it wasn't really used on this particular problem, 
one student was in the habit of giving answers like this if he didn't know how to do a problem.
\end{pfevalenum} 
\end{answer}
\end{eval}

Now we are ready to move on to the main topic of this section.

\begin{df}
\index{logically equivalent}
\index{equivalent!logically}
\index{a logiczzzz@$=$ (logically equivalent)}
Let $p$ and $q$ be propositions.  Then $p$ and $q$ are said to be {\bf logically equivalent}
if $p \leftrightarrow q$ is a tautology.
An alternative (but equivalent) definition is that 
$p$ and $q$ are equivalent if they have the same truth table.  That is,
if they have the same truth value for all assignments of truth values to the variables. 

When $p$ and $q$ are equivalent, we write \boldmath$p=q$.  
An alternative notation is \boldmath$p \equiv q$.
\end{df}

\begin{rem}
$p=q$ is {\bf not} a compound proposition. Rather it is a statement about the relationship between
two propositions.
\end{rem}


There are many {\em logical equivalences} (or {\em identities}/{\em rules}/{\em laws}) that come in handy when working with 
compound propositions.  
Many of them (e.g. commutative, associative, distributive)
will resemble the arithmetic laws you learned in grade school.  Others are very different.
The most common ones are given in Table~\ref{table:logicLaws}.

\begin{table}[tbh]
\centering
$
\begin{array}{|l|l|}
\hline
Name&Equivalence\\
\hline
\hline
commutativity	& p\vee q=q\vee p\\
				& p\wedge q = q\wedge p\\
\hline
associativity	&p\vee(q\vee r)=(p\vee q)\vee r\\
				&p\wedge(q\wedge r)=(p\wedge q)\wedge r\\
\hline
distributive&	p\wedge(q\vee r) = (p\wedge q) \vee (p\wedge r)\\
			&	p\vee(q\wedge r) = (p\vee q)\wedge(p\vee r)\\
\hline
identity&		p \vee F = p\\
			&		p\wedge T = p\\
\hline
negation&p\vee \neg p = T\\
		& p\wedge\neg p = F\\
\hline
domination&	  	 p\vee T \ = T\\
			&	p\wedge F = F\\
\hline
idempotent&		p\vee p=p\\
			&	 p\wedge p = p\\
\hline
double\ negation	& \neg (\neg p)=p	\\
\hline
DeMorgan's	&	 \neg( p \vee q) = \neg p\wedge\neg q\\
			&	\neg (p\wedge q) = \neg p\vee\neg q\\
\hline
absorption & p\vee(p\wedge q) = p\\
& p\wedge(p\vee q) = p\\
\hline
\end{array}
$
\caption{Common Logical Equivalences}\label{table:logicLaws}
\end{table}



We will provide proofs of some of these so you can get the hang of how to prove propositions
are equivalent.  One method is to demonstrate that the propositions have the same truth tables.
That is, they have the same value on every row of the truth table.  But just drawing a truth table
isn't enough.  A statement like ``since $p$ and $q$ have the same truth table, $p=q$'' is necessary
to make a connection between the truth table and the equivalence of the propositions. 
Let's see a few examples.

\begin{exa}
Prove the double negation law: $\neg (\neg p) = p$.
\begin{pf}
The following is the truth table for $p$ and $\neg (\neg p)$.  

\centering
\begin{tabular}{|c|cc|}
\hline $p$ & $\neg p $ & $\neg (\neg p)$ \\
\hline
$T$ & $F $ & $T $ \\
$F$ & $T$ & $F$ \\
\hline
\end{tabular}

Since the entries for both $p$ and $\neg(\neg p)$ are the same for every row, $\neg (\neg p) = p$.
\end{pf}
\end{exa}
  
The two versions of De Morgan's Law
\index{DeMorgan's Law!for propositions} 
are among the most important propositional equivalences.
It is easy to make a mistake when trying to simplify expressions conditional statements, 
and a solid understanding of De Morgan's
Laws goes a long way.  Let's take a look at them.


\begin{exa}
\index{DeMorgan's Law!for propositions}
Prove the first version of DeMorgan's Law: 
 $\neg (p\vee q) = \neg p \wedge \neg q$
\begin{pf}
We can prove this by showing that both expression have the same truth table.
Below is the truth table for $\neg (p\vee q)$ and  $\neg p \wedge \neg q$ (the gray columns).

\centering
\begin{tabular}{|cc|c|g|c|c|g|} \hline $p$ & $q$ & $p\vee q$ &
$\neg (p\vee q)$ & $\neg p$ &
$\neg q$ & $\neg p \wedge \neg q$ \\
\hline 
$T$ & $T$ & $T$ & $F $ & $F $ & $F $ & $F$ \\
$T$ & $F$ & $T$ & $F $ & $F $ & $T $ & $F$ \\
$F$ & $T$ & $T$ & $F $ & $T $ & $F $ & $F$ \\
$F$ & $F$ & $F$ & $T $ & $T $ & $T $ & $T$ \\
\hline
\end{tabular}

Since they are the same for every row of the table,  $\neg (p\vee q) = \neg p \wedge \neg q$.
\end{pf}
\end{exa}

\begin{exer}
Prove the second version of De Morgan's Law: 
$\neg (p\wedge q) = \neg p \vee \neg q$.\\
\pflinetwo
\bigskip

\centering
{\large 
\begin{tabular}{|cc||c|g|c|c|g|} 
\hline 
$p$ & $q$ & $p\wedge q$ &
$\neg (p\wedge q)$ & $\neg p$ &
$\neg q$ & $\neg p \vee \neg q$ \\
\hline 
$T$ & $T$ & & &&&  \\
$T$ & $F$ & & &&& \\
$F$ & $T$ & & &&& \\ 
$F$ & $F$ & & &&& \\
\hline
\end{tabular}
}
\begin{answer}
Below is the truth table for $\neg (p\wedge q)$ and  $\neg p \vee \neg q$ (the gray columns).

\centering
\begin{tabular}{|cc||c|g|c|c|g|} 
\hline 
$p$ & $q$ & $p\wedge q$ & $\neg (p\wedge q)$ & $\neg p$ & $\neg q$ & $\neg p \vee \neg q$ \\
\hline 
$T$ & $T$ & $T$ & $F $ & $F $ & $F $ & $F$ \\
$T$ & $F$ & $F$ & $T $ & $F $ & $T $ & $T$ \\
$F$ & $T$ & $F$ & $T $ & $T $ & $F $ & $T$ \\ 
$F$ & $F$ & $F$ & $T $ & $T $ & $T $ & $T$ \\
\hline
\end{tabular}

Since they are the same for every row of the table,  $\neg (p\wedge q) = \neg p \vee \neg q$.
\end{answer}
\end{exer}

Truth tables aren't the only way to prove that two propositions are equivalent.
You can also use other equivalences.  Let's see an example.

\begin{exa}
Prove the idempotent law
$p \vee p = p$ by using the
other equivalences.   
\begin{pf}
It is easier to prove backwards ($p=p\vee p$). We have 
$$
\begin{array}{llll}
p &=&p \vee F &\mbox{(by identity)}\\
&=& p \vee (p\wedge \neg p) &\mbox{(by negation)}\\
&=&(p \vee p)\wedge(p \vee\neg p) &\mbox{(by distribution)}\\
&=&(p\vee p)\wedge T &\mbox{(by negation)}\\
&=&p\vee p &\mbox{(by identity)}\\
\end{array}
$$
Thus, $p \vee p = p$.
\end{pf}
\end{exa}

\begin{fib}
Prove the idempotent law  $p\wedge p  = p$ by using the
other equivalences.   
\begin{pf}
Notice that
$$
\begin{array}{llll}
p &=&\blankline{2} &\mbox{(by identity)}\\
&=&\blankline{2} &\mbox{(by negation)}\\
&=&\blankline{2} &\mbox{(by distributive)}\\
&=&\blankline{2}&\mbox{(by negation)}\\
&=& p\wedge p &\mbox{(by \blankline{1})}\\
\end{array} $$  
Thus, \blankline{2}.
\end{pf}
\begin{answer}
$$
\begin{array}{llll}
p &=&\underline{p \wedge T} &\mbox{(identity)}\\
&=&\underline{ p \wedge (p\vee \neg p)} &\mbox{(negation)}\\
&=&\underline{(p \wedge p)\vee(p \wedge\neg p)} &\mbox{(distributive)}\\
&=&\underline{(p\wedge p)\vee F}&\mbox{(negation)}\\
&=& p\wedge p &\mbox{(\underline{identity})}\\
\end{array} $$ 
Thus, \underline{$p\wedge p  = p$}.
\end{answer}
\end{fib}

Although it is helpful to specifically state which rules are being used at every step,
it isn't always required.

\begin{exa}\label{thm:thing}
Prove that $(p\wedge q) \vee (p\wedge \neg q) = p$.
\begin{pf} It is not too difficult to see that
 $$(p\wedge q) \vee (p\wedge \neg q) = p\wedge (q \vee \neg q) = p\wedge T = p.$$
\end{pf}
\end{exa}

\begin{exer}
Use the other equivalences (not a truth table) to prove the Absorption laws.
\begin{lenum}
\item
Prove that $p\vee(p\wedge q) = p$.\\
Proof:\\ 
\vspace*{1.75in}
\item
Prove that $p\wedge(p\vee q) = p$.\\
Proof:\\ 
\vspace*{1.75in}
\end{lenum}
\begin{answer}
(a) We can use the identity, distributive, and dominations laws to see that
$$p\vee(p\wedge q) = (p\wedge T)\vee (p\wedge q)=p\wedge(T\vee q) = p\wedge T = p.$$
(b) We can prove this similarly to the previous one, or we can use the previous one
along with distribution and idempotent laws:
$$p\wedge(p\vee q) = (p\wedge p)\vee(p\wedge q) = p\vee(p\wedge q)=p.$$    
\end{answer}
\end{exer}

One use of propositional equivalences is to simplify logical expressions.

\begin{exa} 
Simplify  $\neg(p\vee \neg q)$.
\begin{sol}
 Using DeMorgan's Law and double negation, we can see that 
 $$\neg (p\vee \neg q) = \neg p \wedge \neg (\neg q) = \neg p \wedge q.$$
\end{sol}
\end{exa}
%% Moved to algorithms
%Of course, this also applies to simplifying conditional expressions in code.
%
%\begin{exa}
%Simplify the following code as much as possible. 
%\begin{code} 
%    if ( !(a==null || a.size()<=0) ) {
%        a.clear();
%    }
%\end{code}
%\begin{sol}
%First, notice that by DeMorgan's Law,  \verb+!(a==null || a.size()<=0)+ is equivalent to
%\verb+!(a==null) && !(a.size()<=0)+.  Simplifying a bit more, we get 
%\verb+a!=null && a.size()>0+.  Thus, the code becomes:
%\begin{code}
%    if (a!=null && a.size()>0) {
%        a.clear();
%    }
%\end{code}
%This may not look much simpler, but it is much easier to understand.
%
%This simplification can also be done by defining $p$ to be \verb+a==null+ and $q$ to be \verb+a.size()<=0+.
%Then the expression is $\neg(p\vee q)$.  Applying De Morgan's Law, this is the same as $\neg p\wedge \neg q$,
%which we translate back to \verb+!(a==null) && !(a.size()<=0)+ and simplify as in the final step above.
%\end{sol}
%\end{exa}
%
%As the previous example demonstrates, you can apply the rules to propositions in various forms.
%Sometimes it is useful to explicitly define $p$ and $q$ (and sometimes $r$) and write expressions using
%formal mathematical notation, but at other times it is just as easy to apply the rules the the expressions
%as they are.  In the previous example, we didn't gain that much by defining $p$ and $q$.  
%But with more complicated expressions it certainly can be helpful.
%
%\begin{rem}
%A common mistake is to forget to use De Morgan's law when dealing with negation.
%For instance, in the last example, replacing the code \verb+!(a==null || a.size()<=0 )+
%with the code \verb+!(a==null) || !(a.size()<=0)+ would be incorrect.  
%You cannot just distribute a negation among other terms.  
%Always remember to use De Morgan's law: $\neg(p\vee q) \neq \neg p\vee\neg q$.
%\end{rem}
%
%
%\begin{exer}
%Simplify the following code as much as possible.\\
%Hint: Example~\ref{thm:thing} might be of use.
%\begin{code} 
%    if ((x>0 && x<y) || (x>0 && x>=y)) {
%        x=y;
%    }
%\end{code}
%\vspace*{2in}
%
%\begin{answer}
%Let $p=$``$x>0$'' and $q=$``$x<y$''.  
%Then the conditional above can be expressed as $(p\wedge q)\vee (p\wedge \neg q)$.
%According to Example~\ref{thm:thing}, this is just $p$.  
%Therefore the code simplifies to:
%\begin{code} 
%    if (x>0) {
%        x=y;
%    }
%\end{code}
%\end{answer}
%\end{exer}
%
%
%\begin{eval}
%Simplify the following code as much as possible. 
%\begin{code} 
%    if (x>0) {
%        if(x<y || x>0) {
%            x=y;
%        }
%    }
%\end{code}
%\begin{ssol}
%Because the second if is in the first one which is if ($x>0$) then $x>0$ is duplicated but at
%the same time to satisfy the second one we just need to keep the second if and cut the first one.
%$x<y$ and $x>0$ are independent conditions so they cannot be more simplified.
%So the answer is: 
%\begin{code}
%    if(x<y || x>0) {
%        x=y;
%    }
%\end{code}
%\end{ssol}
%\evallinetwo
%\begin{answer}
%This is not equivalent to the original code.  Consider the case when $x=-1$ and $y=1$, for instance.
%\end{answer}
%\end{eval}
%
%\begin{exer}
%Simplify the following code as much as possible. 
%\begin{code} 
%    if (x>0) {
%        if(x<y || x>0) {
%            x=y;
%        }
%    }
%\end{code}
%\vspace*{1.5in}
%\begin{answer}
%Since both conditions need to be true when if statements are nested, it is the same thing as a conjunction.
%In other words, the two ifs are equivalent to \verb+if( x>0 && (x<y || x>0) )+.
%By absorption, this is equivalent to \verb+if(x>0)+.
%So the simplified code is:
%\begin{code}
%    if(x>0) {
%        x=y;
%    }
%\end{code}
%You can also think about it this way.  The assignment \verb+x=y+ cannot happen unless \verb+x>0+ due to the outer
%if.\footnote{If you think about it, this is why the solution to this in the previous example failed.}
%But the inner if has a disjunction, one part of which is \verb+x>0+, which we already know is true.
%In other words, it doesn't matter whether or not \verb+x<y+.  This argument also leads to the solution we just gave.
%\end{answer}
%\end{exer}
%
%Although some of these examples may seem a bit contrived, in some sense they {\em are} realistic.  As code is
%refactored, code is added and removed in various places, conditionals are combined or separated, etc. and sometimes
%it leads to conditionals that are more complicated than they need to be.  In addition, when working on large teams,
%you will often work on code written by others.  Since some programmers don't have a good grasp on logic, you will
%certainly run into conditional statements that are way more complicated and convoluted than necessary. 
%As I believe these examples demonstrate, simplifying conditionals is not nearly as easy as one might think.
%It takes great care to ensure that your simplified version is still correct.
%
%\begin{rem}
%There is an important difference between the logical operators as discussed here and 
%how they are implemented in programming languages such as Java, C, and C++.
%It is something that is sometimes called {\bf short circuiting}.
%You are probably familiar with the concept even if you haven't heard it called that before.
%It exploits the {\bf domination laws:}
%$$F \wedge q = F$$
%$$T \vee q = T$$
%Let's see an example.
%\end{rem}
%
%\begin{exa}
%Consider the statement \verb+if(x>=0 && a[x]!=0)+.  
%The first domination law implies that when $x<0$, the expression in the if statement will evaluate to 
%false regardless of the truth value of \verb+a[x]!=0+.  Therefore, many languages will simply
%not evaluate the second part of the expression---they will {\em short circuit}.
%
%The same thing happens for statements like \verb+if(x<0 || x>a.length)+. 
%When $x<0$, the expression is true regardless of the truth value of \verb+x>a.length+.  
%Again, many languages don't evaluate the second part of this expression if the first part is true.
%Of course, if the first part is false, the second part is evaluated since the truth value now depends
%on the truth value of the second part.
%
%There are at least two benefits of this.  First, it is more efficient since sometimes less code needs
%to be executed. 
%Second, it allows the checking of one condition before checking a second condition that might cause a crash.
%You have probably used it in statements like the above to make sure you don't index outside the bounds of
%an array.  Another use is to avoid attempting to access methods or fields when a variable refers to
%null (e.g. \verb+if(a!=null && a.size()>0)+).
%
%But this has at least two consequences that can cause subtle problems if you aren't careful.  
%First, although the AND and OR operators are {\em commutative} (e.g. $p\vee q$ and $q\vee p$ are equivalent),
%that is not always the case for Boolean expressions in these languages.  
%For instance, the statement \verb+if(x>=0 && a[x]!=0)+ is not equivalent to \verb+if(a[x]!=0 && x>=0)+ 
%since the second one will cause a crash if $x<0$.
%Second, if the second part of the expression is code that you expect will always be executed,
%you may spend a long time tracking down the bug that this creates.  
%\end{exa}
%
%
%\begin{eval}
%Rewrite the following segment of code so that it is as simple as possible and logically equivalent.
%\begin{code}
%    if( !(list.isEmpty() && list.get(0)>=100) && !(list.get(0)<100) ) {
%        x++;
%    } else {
%        x--;
%    }
%    
%\end{code}
%\begin{ssolenum}
%\item
%The second and third statements mean the same thing.  Also, if the second is true then we got a value so
%we know the list is not empty, so the first statement is unnecessary.  This leads to the following
%equivalent code:
%\begin{code}
%    if(list.get(0) >= 100) {x++;} else {x--;}
%\end{code}
%\evallinetwo
%\item
%I used DeMorgan's law to obtain:
%\begin{code}
%    if(!list.isEmpty() || list.get(0) < 100) { 
%        x++; 
%    } else { 
%        x--; 
%    }
%\end{code}
%\evallinetwo
%\item
%Let $a$ be \verb+list.isEmpty()+ and $b$ be \verb+list.get(0)>=100+.  
%But then $\neg b=$\verb+list.get(0)<100+.
%The original expression is $\neg(a\wedge b)\wedge\neg(\neg b)$.
%But 
%\begin{eqnarray*}
%\neg(a\wedge b)\wedge\neg(\neg b)&=&\neg(a\wedge b)\wedge  b\\
%&=&(\neg a\vee \neg b)\wedge  b\\
%&=&(\neg a\wedge  b)\vee(\neg b\wedge  b)\\
%&=&(\neg a\wedge  b)\vee F\\
%&=&\neg a\wedge  b
%\end{eqnarray*}
%So my simplified code is
%\begin{code}
%    if( !list.isEmpty() && list.get(0)>= 100 ) {
%        x++;
%    } else {
%        x--;
%    }
%\end{code}
%\evallinetwo
%\end{ssolenum}
%\begin{answer}
%\begin{solevalenum}
%\item
%This solution is incorrect.
%There are a few problems.  The obvious one is that the first statement actually prevents the program
%from crashing so it is certainly not unnecessary!  Also, the second and third statements may be equivalent,
%but how are they connected?  For instance, given the expression $\neg (A\wedge B)\vee \neg A$,
%I cannot simply remove the `redundant' $\neg A$ to obtain an ``equivalent'' expression of $\neg (A\wedge B)$
%(if necessary, plug in different truth values for $A$ and $B$ to convince yourself that these are not the same).
%\item
%This is not correct.  The second part of the expression seems to have disappeared.  
%But how can we know it isn't equivalent?  We just need to find a scenario where the two versions do different
%things. Notice that the `simplified' expression is true when the list is not empty regardless of the value of
%element $0$.  But what if the list is not empty and element $0$ is 50?  The original expression is false and the
%`simplified' expression is true.  Clearly not the same.
%\item
%This solution is not only correct, but it is very well argued.
%\end{solevalenum}
%\end{answer}
%\end{eval}
%
%\begin{ques}
%In the solutions to the previous problem we said that the final solution was correct.
%But there might be a catch.  Go back to the original code and the final solution
%and look closer.  Is the final solution
%{\em really} equivalent to the original?  Explain why or why not.\\
%\evallinetwo
%\begin{answer}
%Technically speaking, the final solution is not equivalent.  
%However, it turns out that it is {\em better} than the
%original.  This is because the original code would actually crash if the list is empty.
%Go back and look at the code and verify that this is the case.  Then verify that the 
%final simplified version will {\em not} crash.
%\end{answer}
%\end{ques}
%
%The previous question serves as a reinforcement of a point previously made.  When dealing with logical
%expressions in programs, we have to be careful about our notion of equivalence.  This is because of
%short-circuiting and the fact that expressions in programs, unlike logical statements, can crash instead
%of being true or false.




Table~\ref{tab:otherLaws} contains some important identities involving $\rightarrow$, $\leftrightarrow$, and $\oplus$.
Since these operators are not always present in a programming language,
identities that express them in terms of $\vee$, $\wedge$, and $\neg$ are particularly important.

\begin{table}[h]
\centering
$
\begin{array}{|l|c|l|}
\cline{1-1}\cline{3-3}
p\oplus q = (p\vee q)\wedge \neg(p\wedge q)		&&p\leftrightarrow q = (p\rightarrow q)\wedge (q\rightarrow p)\\
p\oplus q = (p\wedge\neg q)\vee(\neg p\wedge q)	&&p\leftrightarrow q = \neg p\leftrightarrow \neg q\\
\neg(p\oplus q) = p\leftrightarrow q 			&&p\leftrightarrow q = (p\wedge q)\vee (\neg p \wedge \neg q)\\
\cline{1-1}
p\rightarrow q = \neg q \rightarrow \neg p 		&&\neg(p\leftrightarrow q) = p\leftrightarrow\neg q\\
p\rightarrow q = \neg p \vee q					&&\neg(p\leftrightarrow q) = p\oplus q\\
\cline{1-1}\cline{3-3}
\end{array}
$
\caption{Logical equivalences involving $\rightarrow$, $\leftrightarrow$, and $\oplus$}\label{tab:otherLaws}
\end{table}

\begin{exer}
Let $p$ be ``$x>0$'',  $q$ be ``$y>0$,''
and $r$ be ``Exactly one of $x$ or $y$ is greater than 0.''
\begin{enumerate}[(a)]
\item
Express $r$ in terms of $p$ and $q$ using $\oplus$ (and possibly $\neg$).\\
\answerline
\item
Express $r$ in terms of $p$ and $q$ {\em without using $\oplus$}.\\
\answerline
\end{enumerate}
\begin{answer}
(a) $p\oplus q$; 
(b) $(p\wedge \neg q)\vee(\neg p\wedge q)$ or $(p\vee q)\wedge \neg(p\wedge q)$.
Other answers are possible, but most likely you came up with one of these.  
If not, construct a truth table to determine whether or not your answer is correct.
\end{answer}
\end{exer}

Here is the proof of one of the identities from Table~\ref{tab:otherLaws}.

\begin{exa}\label{exa:oplusIdentity}
Prove that $p\oplus q = (p \wedge \neg q) \vee (\neg p \wedge q)$.
\begin{sol}
It is straightforward to see that $(p \wedge \neg q) \vee (\neg p \wedge q)$ is true if
$p$ is true and $q$ is false, or if $p$ is false and $q$ is true, and false otherwise.
Put another way, it is true iff $p$ and $q$ have different truth values. 
But this is the definition of $p\oplus q$.  
Thus, $p\oplus q = (p \wedge \neg q) \vee (\neg p \wedge q)$.
\end{sol}
\end{exa}

The previous example demonstrates an important general principle.
When proving identities (or equations of any sort), sometimes it works best to start from the right hand side.  
More generally, it is often best to start from the more complicated
expression.
Try to keep this in mind in the future.

\newpage

%based on Rosen 1.3 #16
\begin{eval}
Show that $p\leftrightarrow q$ and $(p\wedge q)\vee(\neg p\wedge \neg q)$ are logically equivalent.
\begin{spfenum}
\item
$p\leftrightarrow q$ is true when $p$ and $q$ are both true, and so is $(p\wedge q)\vee(\neg p\wedge \neg q)$.
Therefore they are logically equivalent.

\evallinethree
\item
They are both true when $p$ and $q$ are both true or both false.  Therefore they are logically equivalent.

\evallinethree
\item
Each of these is true precisely when $p$ and $q$ are both true.

\evallinethree
\item
Each of these is true when $p$ and $q$ have the same truth value and false otherwise, so they are
equivalent.

\evallinethree
\end{spfenum}
\begin{answer}
\begin{pfevalenum}
\item
This is an incomplete proof.  It only proves that in one case ($p$ and $q$ both being true) they are
equivalent.  It says nothing about, for instance, whether or not they have the same truth value
when $p$ is true and $q$ is false.
\item
This proof is also incomplete.  It proves that in two cases they have the same truth value, but is silent
about the other cases.  Are we supposed to {\em assume} that in all other cases the expressions are both false?
\item
This is either incomplete or incorrect, depending on how you read it.  
If by ``precisely'' the writer means ``exactly when'', then it is incorrect since the propositions are
also true when both $p$ and $q$ are false.  Otherwise the proof is incomplete because it does not deal
with every case.
\item
This is correct because it exhausts all of the cases.  It is perhaps a bit brief, however.  
The only way I know the proof is actually correct is that I have to verify what the writer said.
By the definition of $p\leftrightarrow q$, what they said is clearly true.  But to see that it is true of
$(p\wedge q)\vee(\neg p\wedge \neg q)$ I have to actually plug in a few values and/or think about the
meaning of the expression.  
\end{pfevalenum}
\end{answer}
\end{eval}

In the previous example, you should have noticed that just a subtle change in wording can be the difference between
a correct or incorrect proof.  When writing proofs, remember to be very precise in how you word things.
You may know what you mean when you wrote something, but a reader can only see what you actually wrote.

\newpage
\section{Predicates and Quantifiers}\label{section:pred}
\begin{df}
\index{predicate}
\index{propositional function}
A {\bf predicate} or {\bf propositional function} is a statement containing one or more variables, 
whose truth or falsity depends on the value(s) assigned to the variable(s).
\end{df}

We have already seen predicates in previous examples.  Let's revisit one.

\begin{exa}
In a previous example we saw that ``$x<0$'' was a contingency, ``$x^2<0$'' was a contradiction, and ``$x^2\geq 0$''
was a tautology.  Each of these is actually a predicate since until we assign a value to $x$,
they are not propositions.  
\end{exa}
Sometimes it is useful to write {\em propositional functions} using functional notation.
\begin{exa}
Let $P(x)$ be ``$x<0$''.  Notice that until we assign some value to $x$, $P(x)$ is neither true nor false.

$P(3)$ is the proposition ``$3<0$,'' which is false.

If we let $Q(x)$ be ``$x^2\geq 0$,'' then $Q(3)$ is ``$3^2\geq 0$,'' which is true.  
\end{exa}

Notice that both $P(x)$ and ``$x<0$'' are propositional functions.  In other words, we don't have to use functional
notation to represent a propositional function.  Any statement that has a variable in it is a propositional function,
whether we label it or not.

\begin{exer}
Which of the following are propositional functions?
\begin{enumerate}[(a)]
\item \blankline{.25} $x^2+2x+1=0$
\item \blankline{.25} $3^2+2\cdot 3 + 1 = 0$
\item \blankline{.25} John Cusack was in movie $M$.
\item \blankline{.25} $x$ is an even integer if and only if $x=2k$ for some integer $k$.
\end{enumerate}
\begin{answer}
(a) is a predicate since it can be true or false depending on the value of $x$.;
(b) is not a predicate since it is simply a false statement--it doesn't contain any variables.;
(c) is a predicate since it can be true or false depending on the value of $M$.;
(d) is not a predicate.  This one is tricky.  This is a definition.
In this statement, $x$ is not a variable but a label for a number so that it can be referred to later in the sentence.
\end{answer}
\end{exer}


\begin{df}
\index{quantifier!universal}
\index{universal quantifier}
\index{for all}
\index{a all@$\forall$ (for all)}
The symbol $\forall$ is the {\bf universal quantifier}, and it is read as ``for all'',
``for each'', ``for every'', etc.
For instance, $\forall x$ means ``for all $x$''.  When it precedes a statement, 
it means that the statement is true {\em  for all values of $x$}.

As the name suggests, the ``all'' refers to everything in the 
{\bf universe of discourse} (or {\bf domain of discourse}, or simply {\bf domain}), which 
is simply the set of objects to which the current discussion relates.
\end{df}

Hopefully you recall that $\naturals$ is the set of natural numbers 
(i.e. $\{0, 1, 2, \ldots\})$,
$\BBZ$ is the set of integers, 
and $\BBZ^+$ is the set of 
positive integers (i.e. $\{1, 2, 3, \ldots\}$).
We will use these in some of the following examples.

\begin{exa}
Let $P(x)$=``$x<0$''.  Then $P(x)$ is a propositional function, and $\forall x P(x)$ means ``all values of $x$ are negative.''
If the domain is $\BBZ$, $\forall x P(x)$ is false.  
However, if the domain is negative integers, $\forall x P(x)$ is true.
\end{exa}


\begin{exa}
Express each of the following English sentences using the universal quantifier.
Don't worry about whether or not the statements are true.
Assume the domain is real numbers.
\begin{enumerate}[(a)]
\item
The square of every number is non-negative.
\item
All numbers are positive.
\end{enumerate}
\begin{sol}
(a) $\forall x (x^2\geq 0)$ \hspace*{.5in}
(b) $\forall x (x >0)$
\end{sol}
\end{exa}

\begin{exer}
Express each of the following using the universal quantifier.
Assume the domain is $\BBZ$.
\begin{enumerate}[(a)]
\item
Two times any number is less than three times that number.\\
\answerline
\item
$n!$ is always less than $n^n$.\\
\answerline
\end{enumerate}
\begin{answer}
\begin{enumerate}[(a)]
\item $\forall x (2x<3x)$.  In case it isn't obvious, there is nothing magical about $x$.  
You could also write your answer as $\forall a (2a<3a)$, for instance.
\item $\forall n (n!<n^n$).
\end{enumerate}
\end{answer}
\end{exer}

\begin{exa}
The expression  $\forall x(x^2\geq 0)$ means ``for all values of $x$, $x^2$ is non-negative''. 
But what constitutes {\em all} values?  In other words, what is the domain?
In this case the most logical possibilities are the integers or real numbers since it seems
to be stating something about numbers (rather than people, for example).
In most situations the context should make it clear what the domain is. 
\end{exa}

\begin{exa}
The expression $\forall x\geq 0, x^3\geq 0$ means ``for all positive values of $x$, $x^3\geq 0$.''
There are several other ways of expressing this idea, but this one is probably the most convenient.  
One alternative would be to restrict the domain to
positive numbers and write it as $\forall x (x^3\geq 0)$.  But sometimes you don't want to 
or can't restrict the domain.  

Another way  to express it is $\forall x(x\geq 0 \rightarrow x^3\geq 0)$.
\end{exa}

\begin{exer}
Use the universal quantifier to express the fact that the square of any integer is not zero as long as the
integer is not zero.\\
\answerline
\begin{answer}
$\forall x\neq 0 (x^2\neq 0)$.
Alternatively, $\forall x (x\neq 0 \rightarrow x^2\neq 0)$.
\end{answer}
\end{exer}


\begin{df}
\index{quantifier!existential}
\index{existential quantifier}
\index{a exists@$\exists$ (there exists)}
\index{exists}
The symbol $\exists$ is the {\bf existential quantifier}, and it is read as ``there exists'',
``there is'', ``for some'', etc.
For instance, $\exists x$ means ``For some $x$''.  When it precedes a statement, 
it means that the statement is true for {\em at least one value of $x$} in the universe.
\end{df}

\vspace*{-.1in}

\begin{exa}
Prove that $\exists x (\sqrt{x}=2)$ is true when the domain is the integers.
\begin{proof}
Notice that when $x=4$, $\sqrt{x}=\sqrt{4}=2$, proving the statement.
\end{proof}
\end{exa}

\vspace*{-.1in}

\begin{exer}
Express the sentence ``Some integers are positive'' using quantifiers.
You may assume the domain of the variable(s) is $\BBZ$.\\
\answerline
\begin{answer}
$\exists x (x>0)$.
\end{answer}
\end{exer}


Sometimes you will see {\em nested quantifiers}.  Let's see a few examples.

\begin{exa}
Use quantifiers to express the sentence ``all positive numbers have a square root,''
where the domain is real numbers.
\begin{sol}We can express this as
$\forall (x > 0) \exists y (\sqrt{x}=y)$.
\end{sol}
\end{exa}

\vspace*{-.1in}

\begin{eval}
Express the sentence ``Some integers are even'' using quantifiers.
You may assume the domain of the variable(s) is the integers.
\begin{ssolenum}
\item
$\exists x(x \mbox{ is even})$.\\
\evallinetwo
\item
$\exists x(x/2\in\BBZ)$.\\
\evallinetwo
\item
$\exists x\exists y (x=2y)$.\\
\evallinetwo
\end{ssolenum}
\begin{answer}
\begin{solevalenum}
\item
While perhaps technically correct, this solution is not very good.  It at least uses a quantifier.
But the fact that it includes the phrase ``is even'' suggests that it could be phrased a bit more `mathematically.'
\item
This solution is pretty good.  It is concise, but expresses the idea with mathematical precision.
Although it doesn't directly appeal to the definition of even, it does use a fact that we all know to be
true of even numbers.
\item
This solution is also good.  It clearly uses the definition of even. It is a bit more complicated
since it uses two quantifiers, but I prefer this one slightly over the second solution.  
But that may be because I didn't come up with the second solution and I refuse to admit that
someone had a better solution than what I thought of (which was this one).
\end{solevalenum}
\end{answer}
\end{eval}

\vspace*{-.1in}

\begin{exa}
Translate $\forall \forall \exists\exists$ into English.
\begin{sol}
It means ``{\em for every upside-down $A$ there exists a backwards $E$}.''
This is a geeky math joke that might make sense if you paid attention
in calculus (assuming you ever took calculus, of course).  
%If you don't get it, don't worry about it.  Move along.
%These aren't the droids you're looking for.
\end{sol}
\end{exa}

\begin{exer}
Express the following statement using quantifiers:
``Every integer can be expressed as the sum of two squares.''
Assume the domain for all three variables (did you catch the hint?) is $\BBZ$.\\
\answerline
\begin{answer}
$\forall x\exists y \exists z (x=y^2+z^2)$.
\end{answer}
\end{exer}

%% Deleted several examples of proving weird specific
%% statements involving quantifiers that don't work in this
%% ordering of the book and are too specific and not motivated
%% anyway.
%\begin{fib}
%Prove or disprove the statement from the previous example.
%\begin{pf}
%The statement is false.  Let $x=3$.
%We need to show that no choice of \blankline{1.5} will yield $y^2+z^2=3$.  
%We can restrict $y$ and $z$ to \blankline{2} since the square of a
%negative integer is the same as the square of its absolute value.
%We will do a proof by cases, considering the possible values of $y$.
%
%$y\neq 0$ since $3$ is not \blankline{2}. 
% 
%If $y=1$, we need \blankline{1.5}, which is impossible.
%
%If $y\geq 2$, $y^2\geq 4$, so we need \blankline{1.5}, \blankline{1.9}.
%
%Since we have \blankline{3} and none of them work, the statement is false.
%\end{pf}
%\begin{answer}
%$y$ and $z$;
%non-negative integers;
%a perfect square;  
%$z^2=2$; 
%$z^2\leq -1$;
%which is also impossible;
%%Why was this here?  Removed 9/10/14.
%%non-negative {\em or} at least $0$;
%exhausted/tried all possible values of $y$.
%\end{answer}
%\end{fib}
%
%
%\begin{exa}\label{exa:allExists}
%Prove or disprove that the following statement is true
%$$
%\forall n\in\BBZ^+\ \exists m\in \BBZ^+\ \big((n+7)^2>49+m\big)
%$$
%\begin{sol}
%We need to prove that given any value of $n\in\BBZ^+$, 
%there is some value $m\in\BBZ^+$ such that $(n+7)^2>49+m$.
%Notice that $(n+7)^2>49+m$ iff $n^2 + 14n > m$ 
%(do a little algebra to see why). 
%So for any given value of $n$, we can choose $m=n^2+14n-1$.
%Since $n\geq 1$, $m=n^2+14n-1\geq 1^2+14\cdot 1-1\geq 1$ 
%so $m$ is a positive integer.
%Therefore the statement is true.
%
%Let's rephrase the proof to make sure it is clear.
%Given any positive integer $n$, we can choose $m=n^2+14n-1$
%and we can see that $49+m=49+n^2+14n-1=(n^2+14n+49)-1=(n+7)^2-1<(n+7)^2$,
%so $(n+7)^2>49+m$. 
%Since $n\geq 1$, $m=n^2+14n-1\geq 1^2+14\cdot 1-1\geq 1$
%so $m$ is a positive integer.
%Thus, we have shown that for any $n\in\BBZ^+$ there is some 
%$m\in\BBZ^+$ such that $(n+7)^2>49+m$. That is, we have just proven that 
%$\forall n\in\BBZ^+\ \exists m\in \BBZ^+\ \big((n+7)^2>49+m\big)$.
%%and the statement is true. However, there is a subtly here that is 
%%important that we did not address: $m$ has to be a natural number. 
%%Since $n>3$, $m = n^2+14n -1>0$, so $m\in\naturals$, and our proof is valid.
%% OLD STUFF
%%First, you need to convince yourself that 
%%if we can always find some value of $m$ based on a given value of $n>3$ such that 
%%$(n+7)^2>49+m$, the statement is true.  
%%Notice that  $(n+7)^2>49+m$ iff $n^2 + 14n > m$.
%%So if we take $m$ to be any number smaller than $n^2+14n$, for instance  
%%$m = n^2+14n -1$, then the statement is true, assuming such a value of $m$
%%is at least 0. Since $n>3$, $m = n^2+14n -1>0, so it is a valid choice for $m$.
%\end{sol}
%\end{exa}
%
%There is a subtly in the proof in Exercise~\ref{exa:allExists} 
%that you might have overlooked.
%The next evaluate exercise will help to see if you caught it.
%
%\begin{eval}\label{eval:exists}
%
%Prove or disprove that the following statement is true
%$$
%\forall n\in\BBN\ \exists m\in \BBN\ \big((n+7)^2>49+m\big)
%$$
%\begin{ssol}
%Let $n\in\BBN$. If we let $m=n^2+14n-1$,
%then $49+m=49+n^2+14n-1=(n^2+14n+49)-1=(n+7)^2-1<(n+7)^2$,
%so $(n+7)^2>49+m$. 
%Thus, we have shown that for any $n\in\BBN$ there is some 
%$m\in\BBN$ such that $(n+7)^2>49+m$. That is, we have just proven that 
%$\forall n\in\BBN\ \exists m\in \BBN\ \big((n+7)^2>49+m\big)$.
%
%\evallinetwo
%\end{ssol}
%\begin{answer}
%Notice that the student seems to have copied part of the proof
%from Example~\ref{exa:allExists}, so it looks pretty convincing.
%Unfortunately, the proof is incorrect. 
%Notice that in this exercise the universe is $\BBN$, not $\BBZ^+$.
%Why does that matter?
%Because in this case we cannot always choose $m=n^2+14n-1$ 
%as suggested in the proof. If $n=0$, 
%this would give $m=-1$ which is not a natural number, 
%so it is not a valid choice. Therefore there is some value of $n$
%(namely, $n=0$) for which the statement is false.
%Since it is not true for all values of $n$, it is false.
%\end{answer}
%\end{eval}
%
%Evaluate~\ref{eval:exists} is a good reminder of why mathematicians are so fussy about the details in proofs. 
%Now it's your turn to give a proper solution
%to that problem. If you have been paying attention, this should be an easy one.
%\begin{exer}
%Prove or disprove that the following statement is true
%$$
%\forall n\in\BBN\ \exists m\in \BBN\ \big((n+7)^2>49+m\big)
%$$
%\vspace*{1.5in}
%\begin{answer}
%It is easy to disprove this. If $n=0$, then 
%$(n+7)^2=49$, so we need to find a natural number such that
%$49>49+m$. This requires we choose $m<0$. 
%But $m$ needs to be a natural number, so $m\geq 0$, a contradiction.
%Thus, for $n=0$, there is no $m\in\naturals$ that works.
%Therefore since the statement is not true for all values of $n$,
%the statement is false.
%\end{answer}
%\end{exer}
%
%\begin{exa}
%Prove or disprove that the following statement is equivalent to the statement from
%the Example~\ref{exa:allExists}.
%$$
%\exists m\in\BBZ^+\ \forall n\in \BBZ^+\ \big((n+7)^2>49+m\big)
%$$
%\begin{sol}
%This is almost the same as the expression from the previous example, but the $\forall n\in \BBZ^+$ and
%$\exists m\in\BBZ^+$ have been reversed.  
%Does that change the meaning?  Let's find out.
% 
%The expression in Example~\ref{exa:allExists} is saying something like 
%``For any positive integer $n$, there is some positive integer $m\ldots$''
%In English, the statement in this example is saying something like ``There exists a positive integer $m$ such that
%for any positive integer $n\ldots$''
%Are these different?  Indeed.  
%The one from Example~\ref{exa:allExists} lets us pick a value of $m$ based on the
%value of $n$. In other words, we can pick a different value of $m$ for each value of $n$.
%The one from this example requires that we pick a value of $m$ that will work for {\em all} values of $n$.
%Can you see how that is saying something different?
%\end{sol}
%\end{exa}
%
%\begin{exa}\label{exa:existsAll}
%Prove or disprove that the following statement is true.
%$$
%\exists m\in\BBZ^+\ \forall n\in \BBZ^+\ \big((n+7)^2>49+m\big)
%$$
%\begin{sol}
%This statement is true.
%Let $m=1$. Then if $n\geq 1$,
%$(n+7)^2\geq (1+7)^2=64>49+m$.
%Thus, the choice of $m=1$ works for all values of $n$, making
%the statement true.
%%We need there to be some value of $m$ such that for any $n>3$, 
%%$n^2+14n>m$ (we worked this out in Example~\ref{exa:allExists}). 
%%Can we find an $m$ such that $m<n^2+14n$ for all values of $n>3$?  Sure.  
%%It should be clear that $m=3^2+14\cdot 3 <n^2+14n$ for all values of $n>3$. 
%\end{sol}
%\end{exa}
%
%Notice that although the statements from Example~\ref{exa:allExists}
%and Example~\ref{exa:existsAll} are not equivalent, they are both true.
%But do not let that confuse you---whether or not two statements are true has nothing to do with whether or not they are saying the same thing. 
%Just because the two statements from these examples sound very similar
%and they are both true, that does not at all imply they are equivalent.
%For instance, $n+n=2n$ and $n\cdot n=n^2$ are both true statements, 
%but they have nothing to do with each other.


\begin{exer}
Find a predicate $P(x,y)$ such that $\forall x \exists y P(x,y)$ and $\exists y\forall x P(x,y)$ have
different true values.  Justify your answer.
(Hint:  Think simple.  Will something like ``$x=y$'' or ``$x<y$'' 
work if you choose the appropriate domain?)\\
Answer:\\
\vspace*{2.5in}
\begin{answer}
You may have a different answer, but here is one possibility based on the hint.
If we let $P(x,y)$ be $x<y$ where the domain for both is the real numbers, then 
 $\forall x \exists y (x<y)$ is true since for any given $x$, we can choose $y=x+1$.
 However, $\exists y\forall x (x<y)$ is false since no matter what value we pick for $y$, $x<y$ is false for
 $x=y+1$.  In other words, it is not true for {\em all} values of $x$.
 As with the previous examples, the difference is that in this case we need to have a single value of $y$
 that works for {\em all} values of $x$.
\end{answer}
\end{exer}

\begin{exa}
Let $P(x)$=``$x<0$''.  Then $\neg\forall x P(x)$ means ``it is not the case that all values of $x$ are negative.''
Put more simply, it means ``some value of $x$ is not negative'', which we can write as $\exists x\neg P(x)$.
\end{exa}

What we saw in the last example actually holds for any propositional function. 
 
\begin{thm}[DeMorgan's Laws for quantifiers]
\index{DeMorgan's Law!for quantifiers}
\index{negation!quantifiers}
\index{a logicn@$\neg$ (NOT)}
If $P(x)$ is a propositional function, then
$$\neg\forall x P(x) = \exists x \neg P(x), \mbox{ and}$$
$$\neg\exists x P(x) = \forall x \neg P(x).$$
\begin{pf}
We will prove the first statement.  The proof of the second is very similar.
Notice that $\neg\forall x P(x)$ is true if and only if $\forall x P(x)$ is false.
$\forall x P(x)$ is false if and only if there is at least one $x$ for which $P(x)$ is false.
This is true if and only if $\neg P(x)$ is true for some $x$.
But this is exactly the same thing as $\exists x \neg P(x)$, proving the result.
\end{pf}
\end{thm}



\begin{exa}\label{exa:allIntsEven}
Negate the following expression, but simplify it so it does not contain the $\neg$ symbol.
$$
\forall n \exists m (2m=n)
$$
\begin{sol}
\begin{eqnarray*}
\neg\forall n \exists m (2m=n)&=&\exists n \neg\exists m (2m=n)\\
&=&\exists n \forall m\neg (2m=n)\\
&=&\exists n \forall m (2m\neq n)\\
\end{eqnarray*}
\end{sol}
\end{exa}

\begin{exer}
Answer the following questions about the expression from Example~\ref{exa:allIntsEven},
assuming the domain is $\BBZ$.
\begin{enumerate}[(a)]
\item
Write the expression in English.  You can start with a direct translation, but then smooth it out as much as possible.\\
\answerlinetwo
\item
Write the negation of the expression in English. State it as simply as possible.\\
\answerlinetwo
\item
What is the truth value of the expression?  Prove it.\\
\answerlinetwo
\end{enumerate}
\begin{answer}
\begin{enumerate}[(a)]
\item
It is saying that every integer can be written as two times another integer.
Simplified, it is saying that every integer is even.
\item
The most direct translation of the final line of the solution is ``There is some integer that cannot be
written as two times another integer for any integer.''  A smoothed-out translation would be ``There is at
least one odd integer.''
\item
Since $3$ is odd, the statement is clearly false.
\end{enumerate}
\end{answer}
\end{exer}

\newpage
%% Moved to algorithms.
%Let's see how quantifiers connect to algorithms.  
%If you want to determine whether or not something (e.g. $P(x)$) is true for all values in a domain
%(e.g., you want to determine the truth value of $\forall x P(x)$), 
%one method is to simply loop through all of the values and test whether or not $P(x)$ is true.  
%If it is false for any value, you know the answer is false.  
%If you test them all and none of them were false, you know it is true.
%
%
%\begin{exa}
%Here is how you might determine if $\forall x P(x)$ is true or false for the 
%domain $\{0, 1, 2, \ldots, 99\}$:
%{\small
%\begin{code}
%    boolean isTrueForAll() {
%        for(int i=0;i<100;i++) {
%            if( !P(i) ) {
%                return false;
%            }
%        }
%        return true;
%    }
%\end{code}
%}
%
%Notice the negation in the code---this can trip you up if you aren't careful.
%\end{exa}
%
%\begin{exa}
%Let $P(x)$ and $Q(x)$ be predicates and the domain be $\{0, 1, 2, \ldots, 99\}$.
%What is \verb+isTrueForAll2()+ determining?  
%\begin{code}
%	    boolean isTrueForAll2() {
%	        for(int i=0;i<100;i++) {
%	            if( !P(i) && !Q(i) )
%	                return false;
%	        }
%	        return true;
%	    }
%\end{code}
%\begin{sol}
%Notice that if both $P(i)$ and $Q(i)$ are false for the same value of $i$, 
%it returns false, and otherwise it returns true.  Put another way, it returns
%true if for every value of $i$, either $P(i)$ or $Q(i)$ is true. 
%Thus, \verb+isTrueForAll2+ is determining the truth value of $\forall i(P(i)\vee Q(i))$.
%\end{sol}
%\end{exa}
%
%\begin{exer}
%Rewrite the expression \verb+( !P(i) && !Q(i) )+ from the previous example so that it uses only one negation.\\
%Answer:\\
%\vspace*{1in}
%\begin{answer}
%Using De Morgan's Law, we get \verb+!(P(i) || Q(i))+.
%\end{answer}
%\end{exer}
%
%\newpage
%\begin{exer}
%Let $P(x)$ and $Q(x)$ be predicates and the domain be $\{0, 1, 2, \ldots, 99\}$.
%What is \verb+isTrueForAll3()+ determining?
%\begin{code}
%		boolean isTrueForAll3() {
%		    boolean result = true;
%		    for(int i=0;i<100;i++) {
%		        if(!P(i)) {
%		            result = false;
%		        }
%		    }
%		    if(result==true) {
%		        return true;
%		    }
%		    for(int i=0;i<100;i++) {
%		        if(!Q(i)) {
%		            return false;
%		        }
%		    }
%		    return true;
%		}
%\end{code}
%\answerlinethree
%\begin{answer}
%First notice that if $P(i)$ is true for every value of $i$, \verb+result+ will be true at the
%end of the first loop, so  \verb+isTrueForAll3+ will return true without even considering $Q$.
%However, if $P(i)$ is false for any value of $i$, then it will go onto the second loop.
%The second loop will return false if $Q(i)$ is false for any value of $i$.  But if $Q(i)$ is true
%for all values of $i$, the method returns true.
%So, how do we put this all together into a simple answer?
%Notice that the only time it returns true is if either $P(i)$ is always true or if $Q(i)$ is always true.
%In other words,  \verb+isTrueForAll3+ is determining the truth value of $(\forall i P(i)) \vee (\forall i Q(i))$.
%By the way, notice that I used the variable \verb+i+ for both quantifiers. This is sort of like using
%the same variables for two separate loops.
%\end{answer}
%\end{exer}
%
%\begin{exa}
%Now we are ready for the million dollar 
%question:\footnote{There is no million dollars for answering this question.  It's just an expression.}
% Are \verb+isTrueForAll2+ and \verb+isTrueForAll3+ determining the same thing?
%\begin{sol}
%At first glance, it looks like they might be.  But we need to dig deeper, and we need to prove one way or the other.
%To prove it, we would need to show that these expressions evaluate to the same truth value, regardless of what $P$ and $Q$ are.   
%To disprove it, we just need to find a $P$ and a $Q$ for which these expressions have different truth values.
%But let's first talk it through to see if we can figure out which answer seems to be correct.
%
%$\forall i(P(i)\vee Q(i))$ is saying that for every value of $i$, either $P(i)$ or $Q(i)$ has to be true.
%$\forall i P(i) \vee \forall i Q(i)$ is saying that either $P(i)$ has to be true for every $i$, or that
%$Q(i)$ has to be true for every $i$.  These sound similar, but not exactly the same, so we cannot be sure yet.
%In particular, we cannot jump to the conclusion that they are not equivalent because we described each with
%different words.  There are many ways of wording the same concept. 
%
%At this point we either need to try to tweak the wording so that we can see that they are really saying the same
%thing, or we need to try to convince ourselves they aren't.  
%Let's try the latter.
%
%What if $P(i)$ is always true and $Q(i)$ is always false?  
%Then  both statements are true.  But that doesn't prove that they are always both true, so this doesn't help.
%
%Let's try something else.  
%What if we can find a $P(i)$ and a $Q(i)$ such that for any given value of $i$, we
%can ensure that either $P(i)$ or $Q(i)$ is true, but also that there is some value $j$ such that $P(j)$ is false
%and some value $k\neq j$ such that $Q(k)$ is false?  
%Then $\forall i(P(i)\vee Q(i))$ would be true, but $\forall i P(i) \vee \forall i Q(i)$ false, so this would work. 
%But in order to be certain, 
%we have to know that such a $P$ and $Q$ exist.\footnote{Consider this:  If I can find an even number
%that is prime but is not $2$, then there would be at least 2 even primes.  That's great.  Unfortunately, 
%I can't find such a number.}
%
%What if we let $P(i)$ be ``$i$ is even'', $Q(i)$ be ``$i$ is odd'', and
%the universe be $\BBZ$.  
%Then $\forall i P(i) = \forall i Q(i) = F$, so $\forall i P(i) \vee \forall i Q(i)=F$, 
%but  $\forall i(P(i)\vee Q(i))=T$. 
%Now we have all of the pieces.  Let's put this all together in the form of a proof.
%
%\begin{pf} (that  $\forall i(P(i)\vee Q(i)) \neq \forall i P(i) \vee \forall i Q(i)$)\\
%Let $P(i)$ be ``$i$ is even'', $Q(i)$ be ``$i$ is odd'', and
%the universe be $\BBZ$.  Then  $\forall i(P(i)\vee Q(i))$ is true since every integer is either even or odd.
%On the other hand, $\forall i P(i)$ is false since there are integers that are not even 
%and $\forall i Q(i)$ is false since there are integers that are not odd.  Thus, 
%$\forall i P(i) \vee \forall i Q(i)$ is false.  Since they have different truth values, 
%  $\forall i(P(i)\vee Q(i)) \neq \forall i P(i) \vee \forall i Q(i)$
%\end{pf}
%\end{sol}
%\end{exa}

%\input{logicSetsBooleanAlg/logicExercises.tex}


%%%%%%%%%%%%%%%%%%%%%%%%%%%%%%%%%%%
%%%%%%%%%%%%%%%%%%%%%%%%%%%%%%%%%%%
%%%%%%%%%%%%%%%%%%%%%%%%%%%%%%%%%%%
%%%%%%%%%%%%%%%%%%%%%%%%%%%%%%%%%%%
%%%%%%%%%%%%%%%%%%%%%%%%%%%%%%%%%%%

\newpage
\section{Normal Forms}\label{section:normal}

Earlier we saw identities that express logical operators
in terms of $\vee$, $\wedge$, and $\neg$.  
It turns out that even if there isn't an identity that does it,
there is a straightforward technique to convert any logical expression
into one only using $\vee$, $\wedge$, and $\neg$.  
That is the topic of this section. 

Specifically, we will introduce two standard forms that every boolean expression can be written in:
{\em disjunctive normal form} and {\em conjunctive normal form}.
These form have connections to important areas of computer science including
circuit design and minimization, artificial intelligence algorithms,
automated theorem proving, and the study of algorithm complexity.
We begin with a few necessary definitions.

\begin{df}\index{literal}
A {\bf literal} is a boolean variable or its negation.
\end{df}

\begin{exa}
Let $p$, $q$, and $r$ be boolean variables.  Then 
$p$, $\neg p$, $q$, $\neg q$, $r$, and $\neg r$ are all literals.

On the other hand, $p\wedge q$, $\neg p \rightarrow q$, and $p\wedge q\wedge r$ 
are {\em not} literals because they include boolean operations of two or more variables.
In other words, none of them are a variable or the negation of a variable.
\end{exa}

\begin{exer}
Let $p$, $q$, and $r$ be boolean variables. 
Which of the following are literals?
$q\vee r$, $\neg p$, $q$, $p\wedge q \wedge r$, $\neg p \wedge q$, $r$.\\
\answerlinetwo 
\begin{answer}
$\neg p$, $q$, and $r$. 
\end{answer}
\end{exer}

\begin{df}\index{conjunctive clause}
A {\bf conjunctive clause} is a conjunction ({\em AND}) of one or more literals.
\end{df}


\begin{exa}
Let $p$, $q$, and $r$ be boolean variables.  
Then 
$p\wedge q\wedge r$, $\neg p\wedge r$, and $r\wedge \neg q \wedge p$ are all conjunctive clauses.


$\neg(p\wedge q)$ is not a conjunctive clause because it has a negation that is applied to the conjunction
and not to just a variable. Other examples that are {\em not} conjunctive clauses are
$p\vee r$, $p\vee q \wedge r$, $ p\leftrightarrow r$, and $p\wedge q \wedge (r\oplus p)$.
\end{exa}


\begin{exa}
Literals are conjunctive clauses since they are a conjunction of a single variable.
This might sound weird because if you only have a single variable, there is nothing to ``conjoin'' it to. 
But it is just like if someone asked to to add up all of the money in your pocket--if you only have a single 
dollar, you will say you have one dollar, having ``added up'' the dollar.

Therefore, 
$p$, $\neg p$, $q$, $\neg q$, $r$, and $\neg r$ are all conjunctive clauses.
\end{exa}

\begin{exer}
Let $p$, $q$, and $r$ be boolean variables. 
Which of the following are conjunctive clauses?\\
$q\vee r$, 
$\neg p$,
$q$, 
$p\wedge q \wedge r$, 
$\neg p \rightarrow q$, 
$\neg p \wedge q$, 
$r$, 
$\neg r \wedge p \wedge q$,
$q\vee\neg r$, 
$p \wedge \neg(r\wedge q)$
\\
\answerlinetwo
\begin{answer}
$\neg p$,
$q$, 
$p\wedge q \wedge r$, 
$\neg p \wedge q$, 
$r$, 
$\neg r \wedge p \wedge q$.
\end{answer}
\end{exer}


\begin{df}\index{disjunctive normal form}\index{sum-of-products}
A logical expression is in {\bf disjunctive normal form} ({\bf DNF}) 
(or {\bf sum-of-products expansion})
if it is expressed as a disjunction ({\em OR}) of conjunctive clauses.
\end{df}

\begin{exa}
Let $p$, $q$, and $r$ be boolean variables.
Then the following are in disjunctive normal form:
\begin{itemize}
\item 
$q$ (It is the disjunction of a single conjunctive clause that consists of just a literal.)
\item
$p\wedge \neg q$  (It is the disjunction of a single conjunctive clause.)
\item
$p\vee \neg q$ (It is the disjunction of two conjunctive clauses, each of which is just a literal.)
\item
$(p\wedge q\wedge r)\vee(\neg p\wedge r)$
\item
$p\vee(q\wedge \neg p)\vee(r\wedge\neg p)$
\item
$r\wedge \neg q \wedge p$
\end{itemize}

\noindent
These are {\em not} in disjunctive normal form.
\begin{itemize}
\item
$p\rightarrow q$
\item
$p\wedge (q \vee r)$
\item
$p\vee \neg (r \wedge q)$
\item
$p\vee(q\wedge \neg p)\wedge(r\vee \neg q)$
\item
$(p\leftrightarrow q)\vee (q\wedge \neg r)\vee \neg p$
\end{itemize}
\end{exa}


\begin{exer}
Let $p$, $q$, and $r$ be boolean variables. 
Which of the following are in disjunctive normal form?\\
$\neg p$,
$q\vee r$, 
$\neg q\wedge r$, 
$p\wedge q \wedge r$, 
$(\neg p \rightarrow q)\vee (q\wedge r)$,  
$\neg(p \wedge \neg q)$,
$\neg r \wedge p \wedge q$,
$\neg(p \vee q)$,
$p \wedge \neg(r\wedge q)$,
$(p\wedge \neg r)\vee(r\wedge q) \vee (\neg q\wedge p)$,
$(p\vee \neg r)\wedge(r\vee q) \wedge (\neg q\vee p)$,
$(p\wedge \neg r)\vee(r\vee q) \vee (\neg q\wedge p)$,
$(p\wedge \neg r)\vee(r\wedge q) \wedge (\neg q\wedge p)$,
$(p\wedge \neg r)\vee(r\wedge q) \vee (\neg q\wedge p\wedge r)$
\\
\answerlinetwo
\begin{answer}
$\neg p$,
$q\vee r$, 
$\neg q\wedge r$, 
$p\wedge q \wedge r$,
$\neg r \wedge p \wedge q$,
$(p\wedge \neg r)\vee(r\wedge q) \vee (\neg q\wedge p)$,
$(p\wedge \neg r)\vee(r\wedge q) \vee (\neg q\wedge p\wedge r)$
\end{answer}
\end{exer}

Make sure you have a clear understanding of when an expression is and is not a literal, 
a conjunctive clause, or in disjunctive normal form.

Now you understand what disjunctive normal form is and can recognize when an expression is 
in this form. Next we describe how to convert any expression to an equivalent one
that is in disjunctive normal form. The procedure involves constructing a truth table
for the expression. The process is pretty straightforward once you get the hang of it, but
it can be a little tricky at first so pay careful attention!

\begin{proc}\label{dnfConvert}
This will convert a boolean expression to disjunctive normal form.
\begin{enumerate}
\item Create a truth table for the expression.
\item Identify the rows having output  $T$. 
\item For each such row, create a conjunctive clause that includes all of the
variables which are true on that row and the negation of all of the variables that
are false.
\item Combine all of the conjunctive clauses by disjunctions.
\end{enumerate}
\end{proc}

\begin{exa}
Express $p\oplus q$ in disjunctive normal form.
\begin{sol}
The truth table for $p\oplus q$ is given to the right.
The second and third rows of the table are true, so we use those to construct the 
disjunctive normal form.

\begin{minipage}{3.5in}
The second row yields conjunctive clause $p\wedge\neg q$, and the third row
yields conjunctive clause $\neg p\wedge q$.  The disjunction of these is
$(p\wedge\neg q)\vee(\neg p\wedge q)$.

Thus, $p\oplus q = (p\wedge\neg q)\vee(\neg p\wedge q)$.
\end{minipage}
\hfill
\begin{minipage}{1.25in}
\begin{tabular}{|cc||c|}
\hline 
$p$ & $q$ &$p\oplus q$\\
\hline 
$T$ & $T$& $F$ \\
$T$ & $F$& $T$ \\
$F$ & $T$& $T$ \\
$F$ & $F$& $F$ \\
\hline
\end{tabular}
\end{minipage}
\end{sol}
\end{exa}

The previous example is essentially just another proof of the identity that was
proven in Example~\ref{exa:oplusIdentity}.

\begin{exer}
Express $p\leftrightarrow q$ in disjunctive normal form.\\
\begin{minipage}{3in}
\end{minipage}
\hfill
\begin{minipage}{1.5in}
{\Large
\begin{tabular}{|cc||c|}
\hline 
$p$ & $q$ &$p\leftrightarrow q$\\
\hline 
$T$ & $T$& \\
$T$ & $F$& \\
$F$ & $T$& \\
$F$ & $F$& \\
\hline
\end{tabular}
}
\end{minipage}
\vspace*{.5in}

\begin{answer}
\begin{minipage}[t]{4in}
The truth table for $p\leftrightarrow q$ is given to the right.
The first row yields conjunctive clause $p\wedge q$, and the fourth row
yields conjunctive clause $\neg p\wedge \neg q$.  The disjunction of these is
$(p\wedge q)\vee(\neg p\wedge \neg q)$. 
Thus, $p\leftrightarrow q = (p\wedge q)\vee(\neg p\wedge \neg q)$.
\end{minipage}
\hfill
\begin{minipage}[t]{1.5in}
\begin{tabular}[t]{|cc||c|}
\hline 
$p$ & $q$ &$p\leftrightarrow q$\\
\hline 
$T$ & $T$& $T$ \\
$T$ & $F$& $F$ \\
$F$ & $T$& $F$ \\
$F$ & $F$& $T$ \\
\hline
\end{tabular}
\end{minipage}
\end{answer}
\end{exer}


\begin{exa}
Express $Z$ in disjunctive normal form.
$$\begin{array}{|lll|l|}
\hline
p & q & r & Z \\
\hline
T & T & T & T \\
T & T & F & T \\
T & F & T & F \\
T & F & F & F \\
F & T & T & F \\
F & T & F & T \\
F & F & T & F \\
F & F & F & T \\
\hline
\end{array}$$
\begin{sol}
$Z=(p\wedge q \wedge r)\vee
(p\wedge q \wedge \neg r)\vee
(\neg p\wedge q \wedge \neg r)\vee
(\neg p\wedge \neg q \wedge \neg r)$.
\end{sol}
\end{exa}

The solution from the previous example can be simplified to 
$Z=(p\wedge q)\vee (\neg p\wedge \neg r)$.
Although this can be done by applying the logical equivalences
we learned about earlier, there are more sophisticated techniques
that can be used to simplify expressions that are in disjunctive normal form.
This is beyond our scope, but you may learn more about this if you
take a computer organization class and discuss circuit minimization.
The important point I want to make here is that computing the disjunctive normal form
of an expression using the technique we describe will not always produce the
most simple form of the expression.  In fact, most of the time it won't be.

\begin{exer}
Express $Y$ in disjunctive normal form.

$\begin{array}{|lll|l|}
\hline
p & q & r & Y \\
\hline
T & T & T & F \\
T & T & F & T \\
T & F & T & F \\
T & F & F & F \\
F & T & T & T \\
F & T & F & F \\
F & F & T & T \\
F & F & F & T \\
\hline
\end{array}$
\begin{answer} 
$Y=
(p\wedge q \wedge \neg r) \vee
(\neg p\wedge q \wedge r) \vee
(\neg p\wedge \neg q \wedge r) \vee
(\neg p\wedge \neg q \wedge \neg r)$.
\end{answer}
\end{exer}

There is another important form that is very similar to disjunctive normal form.

\begin{df}\index{conjunctive normal form}\index{product-of-sums}
\index{disjunctive clause}
A {\bf disjunctive clause} is a disjunction ({\em OR}) of one or more literals.
A logical expression is in {\bf conjunctive normal form} ({\bf CNF})
(or {\bf product-of-sums expansion})
if it is expressed as a conjunction ({\em AND}) of disjunctive clauses.
\end{df}

There are several methods for converting to conjunctive normal form.
They generally involve using double negation, distributive, and De Morgan's laws
either based on the truth table or based on the disjunctive normal form.
However, we won't discuss these techniques here.  

%%%%%%%%%%%%%%%%%%%%%%%%%%%%%%%%%%%%%%%%%%%%%%%%
